\section{The global theory: models, finite descent, ample envelopes, Siu division and the main theorem}\label{sec:global}

This chapter is the route for Sections~2, 4, 6, 7 and~8 of \cite{FHX} (cited in the revision of 2026-10-28;
outside Section~5 the earlier revisions \texttt{832e086} and \texttt{d64981e} agree verbatim), for the Main Theorem
(\texttt{thm:intro-main}) and for the characterization \texttt{cor:intro-characterization}. The appendix
(singular Yau--Tian--Donaldson) is out of scope. Nothing below depends on \cite[Theorem~6.21]{BJ22}
(\texttt{thm:pshdiv}) or on its consequences (\cite[Corollaries~6.22, 6.41, Theorems~6.31, 6.35, 6.40]{BJ22}).

\begin{itemize}
\item \emph{Model} (\S\ref{subsec:gl-model}). A projective variety with $r$ globally generated line bundles
$L_1,\dots,L_r$ whose product map is an embedding is modelled by a finitely generated $\N^r$-graded $k$-algebra
$S\subseteq K[T_1,\dots,T_r]$ ($K$ the function field) with coordinates in the degrees $e_c$; the
classes are $\theta\in\R^r$, read as $\sum_c\theta_c\,c_1(L_c)\in N^1(X)_\R$. Fubini--Study functions use
the \emph{sections} $R_a$, $a\in\mathbf Z^r$ (rational functions of degree $a$ without poles), so that classes
with negative coordinates are available. The \emph{model-nef cone} $N(\mathcal X)\subseteq\R^r$ is the closure of
the cone of degrees carrying base point free sections; it is intrinsic to the model, need not be simplicial,
and corresponds to the nef cone of $X$ (Remark~\ref{rmk:gl-model-comparison}). The psh theory of
\S\ref{subsec:pg-fs-systems} extends verbatim to classes in $N(\mathcal X)$.
\item \emph{No non-negligibility.} The direction ``envelope property $\Rightarrow$ convergence on $X^{\divv}$'' of
\cite[Proposition~2.4]{FHX}, which needs \cite[Theorem~5.6]{BJ22} (divisorial points are nonnegligible) and
\cite[Corollary~4.17]{BJ22} (finiteness on $X^{\divv}$), is never used: every step of \cite{FHX} that uses it is
rewritten with usc regularizations of suprema, the divisorial comparison theorem, and the fact that usc
regularization is local on open sets (Lemma~\ref{lem:gl-usc-local}).
\item \emph{No energy theory, no test configurations.} Regularization by nets (built into the definition) replaces
\cite[Corollary~12.18]{BJ22}; Step~2 of \cite[Theorem~6.1]{FHX} uses Fubini--Study envelopes and extension of
scalars of divisorial valuation rings instead of base ideals of test configurations.
\item \emph{Proved alterations and remaining inputs.} De Jong's scheme-theoretic theorem and
its quotient-model form (Theorem~\ref{thm:gl-dejong}) are proved. The root public statements
retain local positivity and the characteristic-zero smooth-envelope theorem as hypotheses; local
positivity is proved over every algebraically closed field (Theorem~\ref{thm:lpn-local-positivity-instance}),
and Theorem~\ref{thm:lpn-main-unconditional} states the main theorems without it. The global
compatibility statements also accept fibre constancy explicitly; the root statements supply
it from the proved fibre-constancy theorem. Frobenius descent and additive powers used here
are proved tools in the model.
\item \emph{Classes.} The main statements cover all model classes $\theta\in\R^r$.
Chapter~\ref{sec:nonnef} handles the non-nef step through homogeneous decreasing limits and
structured diagonal lemmas, without an extra PSH2 hypothesis.
\end{itemize}

\subsection{The model}\label{subsec:gl-model}

\begin{convention}\label{conv:gl-setting}
\lean{EnvelopeProperty.MultiModel.piece, EnvelopeProperty.MultiModel.mem_piece, EnvelopeProperty.MultiModel.ofPiece, EnvelopeProperty.MultiModel.homComp, EnvelopeProperty.MultiModel.coordElem, EnvelopeProperty.MultiModel.coordAlg, EnvelopeProperty.MultiModel.monSet, EnvelopeProperty.MultiModel.zMon}
\leanok
\uses{conv:pg-setting, def:berk}
$k$ is algebraically closed with the trivial absolute value. $r\ge1$; $e_1,\dots,e_r$ is the standard basis of
$\mathbf Z^r$, $\mathbf 1=\sum_ce_c$; vectors in $\mathbf Z^r,\Q^r,\R^r$ are ordered coordinatewise, $a\vee b$ is the
coordinatewise maximum, and $a$ is \emph{positive} if all its coordinates are $>0$. For a field $K\supseteq k$,
$K[T]=K[T_1,\dots,T_r]$ and $K[T^{\pm}]$ is the Laurent polynomial ring; $T^a=\prod_cT_c^{a_c}$; an element $fT^a$ ($f\in K$)
is \emph{homogeneous of degree $a$}, and $p_a\in K$ denotes the coefficient of $T^a$ in $p$. A subalgebra of $K[T]$ is
\emph{graded} if it contains the homogeneous components of its elements. Functions take values in $[-\infty,\infty)$;
$f^\star$ is the usc regularization; nets are indexed by directed sets.
\end{convention}

\begin{definition}[Models]\label{def:gl-model}
\lean{EnvelopeProperty.MultiModel, EnvelopeProperty.MultiModel.IsNormal, EnvelopeProperty.MultiModel.isNormal_iff, EnvelopeProperty.MultiModel.isFractionRing_S, EnvelopeProperty.MultiModel.fg_top_intermediateField}
\leanok
\uses{conv:gl-setting, conv:wu-graded}
\source{fhx-envelope-20261028:subsec:conventions, fhx-envelope-20261028:subsec:prelim-berkovich}
A \emph{model} of rank $r$ is a triple $\mathcal X=(K,S,z)$ consisting of a finitely generated field extension
$K\supseteq k$, a finitely generated graded $k$-subalgebra $S\subseteq K[T]$, and \emph{coordinates}
$z^{(c)}_j=f^{(c)}_jT^{e_c}\in S$ ($1\le c\le r$, $0\le j\le M_c$, $f^{(c)}_j\in K$), such that
\begin{enumerate}
\item[(M1)] for every $c$ some $f^{(c)}_j$ is nonzero;
\item[(M2)] $S$ is integral over the subalgebra $k[z]$ generated by the coordinates;
\item[(M3)] $K$ is generated as a field by the quotients $f/g$ with $f,g\in S_a\coloneqq\{h\in K:hT^a\in S\}$,
$a\in\N^r$, $g\ne0$.
\end{enumerate}
$\mathcal X$ is \emph{normal} if $S$ is integrally closed. For $b\in\N^r$, $\mathrm{Mon}(b)$ is the set of multi-indices
$\beta=(\beta^{(c)}_j)$ with $\sum_j\beta^{(c)}_j=b_c$ for all $c$; $z^\beta=\prod(z^{(c)}_j)^{\beta^{(c)}_j}=\tilde z^\beta T^b$ with
$\tilde z^\beta\in K$.
\end{definition}

\begin{remark}[Examples]\label{rmk:gl-model-examples}
\lean{EnvelopeProperty.MultiModel.quotModel, EnvelopeProperty.MultiModel.quotModelEquiv, EnvelopeProperty.MultiModel.ker_quotMap, EnvelopeProperty.MultiModel.normalizationIn, EnvelopeProperty.MultiModel.isNormal_normalizationIn, EnvelopeProperty.MultiModel.normalization, EnvelopeProperty.MultiModel.mem_normalization_S, EnvelopeProperty.MultiModel.finite_normalization, EnvelopeProperty.MultiModel.isGraded_normalizationInS}
\leanok
\uses{def:gl-model, conv:pg-setting, def:pg-blowup, lem:wu-graded-closure, thm:bj-normalization-finite}
(1) \emph{Quotient models.} Let $R=k[x^{(c)}_j]$ be multigraded by $\deg x^{(c)}_j=e_c$ and $I\subseteq R$ a multihomogeneous
prime ideal that is \emph{relevant} (for every $c$ some $x^{(c)}_{j_c}\notin I$). With $K$ the field of degree-zero fractions
of $R/I$, the map $x^{(c)}_j\mapsto(\bar x^{(c)}_j/\bar x^{(c)}_{j_c})T^{e_c}$ has kernel $I$ (a multihomogeneous element
$F$ of degree $a$ maps to $(\bar F/\prod_c(\bar x^{(c)}_{j_c})^{a_c})T^a$, and components of distinct degrees go to distinct
monomials), so its image is a model $\mathcal X_I$ with $S\cong R/I$; $X=V(I)\subseteq\PP^{M_1}\times\dots\times\PP^{M_r}$. For
$r=1$ this is $V(I)\subseteq\PP^M$ of Convention~\ref{conv:pg-setting} (normal iff projectively normal), and for $r=2$ and
$I=I'$ the blowup of Definition~\ref{def:pg-blowup}.
(2) \emph{Normalization.} The integral closure $S^\nu$ of $S$ in $K[T]$ (equal to its integral closure in
$\operatorname{Frac}S=K(T)$, as $K[T]$ is normal) is graded (torus argument of Lemma~\ref{lem:wu-graded-closure}(1)) and finite over
$S$ (Theorem~\ref{thm:bj-normalization-finite}); $\mathcal X^\nu=(K,S^\nu,z)$ is a normal model, the normalization.
(3) \emph{Finite covers.} For a finite extension $E\supseteq K$, the integral closure $S_E$ of $S$ in $E[T]$ is graded and
finite over $S$, and $(E,S_E,z)$ is a normal model, the normalization of $X$ in $E$.
\end{remark}

\begin{lemma}[Multigraded extension]\label{lem:gl-multigraded-extension}
\lean{EnvelopeProperty.MultiModel.IsMultigradedNorm, EnvelopeProperty.MultiModel.multigradedFun, EnvelopeProperty.MultiModel.IsMultigradedNorm.valuation, EnvelopeProperty.MultiModel.IsMultigradedNorm.exists_dominant, EnvelopeProperty.MultiModel.ofMultigradedNorm, EnvelopeProperty.MultiModel.exists_multigradedExtension, EnvelopeProperty.MultiModel.ofMultigradedNorm_homNorm}
\leanok
\uses{def:gl-model, lem:bj-graded-extension, lem:pg-bigraded-extension}
Let $N$ be a function on the homogeneous elements of $S$ with values in $\R_{\ge0}$ such that $N(0)=0$, $N(\lambda)=1$ for
$\lambda\in k^\times$, $N(pq)=N(p)N(q)$, $N(p+q)\le\max(N(p),N(q))$ for $p,q$ of the same degree, and $\max_jN(z^{(c)}_j)=1$
for every $c$. Then $|p|\coloneqq\max_aN(p_aT^a)$ is a multiplicative seminorm on $S$, trivial on $k$, with
$|p|=\max_a|p_aT^a|$ and $\max_j|z^{(c)}_j|=1$.
\end{lemma}

\begin{proof}
Order $\N^r$ lexicographically; this total order is compatible with addition. All properties except multiplicativity are
immediate. Since $(pq)_\gamma T^\gamma=\sum_{\alpha+\beta=\gamma}p_\alpha q_\beta T^\gamma$, $N((pq)_\gamma T^\gamma)\le|p||q|$, which settles
$|p||q|=0$. Otherwise let $\alpha_0$ (resp.\ $\beta_0$) be the largest degree with $N(p_{\alpha_0}T^{\alpha_0})=|p|$ (resp.\
$N(q_{\beta_0}T^{\beta_0})=|q|$). In $(pq)_{\alpha_0+\beta_0}T^{\alpha_0+\beta_0}$ every term other than $p_{\alpha_0}q_{\beta_0}T^{\alpha_0+\beta_0}$ has
$\alpha>\alpha_0$ or $\beta>\beta_0$, hence norm $<|p||q|$, so $N((pq)_{\alpha_0+\beta_0}T^{\alpha_0+\beta_0})=|p||q|$ and $|pq|=|p||q|$.
\end{proof}

\begin{definition}[Points]\label{def:gl-points}
\lean{EnvelopeProperty.MultiModel.IsMultigraded, EnvelopeProperty.MultiModel.IsNormalized, EnvelopeProperty.MultiModel.an, EnvelopeProperty.MultiModel.triv, EnvelopeProperty.MultiModel.HasGenericSupport, EnvelopeProperty.MultiModel.valSet, EnvelopeProperty.MultiModel.fieldVal, EnvelopeProperty.MultiModel.fieldVal_div, EnvelopeProperty.MultiModel.eq_fieldVal}
\leanok
\uses{def:gl-model, def:berk}
\source{fhx-envelope-20261028:subsec:prelim-berkovich, bj22-trivval:sec:Berk}
$\mathcal X^\an$ is the set of multiplicative seminorms $|\cdot|_v$ on $S$, trivial on $k$, that are \emph{multigraded},
$|p|_v=\max_a|p_aT^a|_v$, and \emph{normalized}, $\max_j|z^{(c)}_j|_v=1$ for every $c$, with the topology induced from
$(\Spec S)^\an$ (Definition~\ref{def:berk}). The \emph{trivial point} $\star$ has $|p|_\star=1$ for $p\ne0$. A point has
\emph{generic support} if $|p|_v\ne0$ for all $p\ne0$; these points form $\mathcal X^{\mathrm{val}}$, and such a $v$ defines
an absolute value $|f/g|_v=|fT^a|_v/|gT^a|_v$ on $K$, $f,g\in S_a$. This is well defined because every element of $K$ is such a
single quotient ($K=\operatorname{Frac}A_j$ for any chart, Lemma~\ref{lem:gl-charts}(1)); in Lean it is obtained by extending $\iota_j(v)$ from $A_j$
to $K$ and characterized by this formula.
\end{definition}

\begin{lemma}[Basic properties]\label{lem:gl-points-basic}
\lean{EnvelopeProperty.MultiModel.IsNormalized.val_le_one, EnvelopeProperty.MultiModel.isCompact_an, EnvelopeProperty.MultiModel.instCompactSpaceAn, EnvelopeProperty.MultiModel.instT2SpaceAn, EnvelopeProperty.MultiModel.instSecondCountableTopologyAn, EnvelopeProperty.MultiModel.IsNormalized.exists_val_zMon_eq_one, EnvelopeProperty.MultiModel.IsNormalized.sup_monSet_val_zMon, EnvelopeProperty.MultiModel.continuous_iff_an}
\leanok
\uses{def:gl-points, lem:polydisc-compact, lem:gl-multigraded-extension}
\begin{enumerate}
\item $|p|_v\le1$ for all $p\in S$ and $v\in\mathcal X^\an$.
\item $\mathcal X^\an$ is compact Hausdorff, $\star\in\mathcal X^\an$, and $\mathcal X^\an$ is second countable if $k$ is countable.
\end{enumerate}
\end{lemma}

\begin{proof}
(1) It suffices to treat homogeneous $p$ of degree $a$. Since $S$ is graded and integral over the graded subalgebra $k[z]$,
taking homogeneous components of an integral equation gives $p^n+\sum_{i=1}^nc_ip^{n-i}=0$ with $c_i\in k[z]$ homogeneous of degree
$ia$ (as in Lemma~\ref{lem:wu-graded-closure}). Each $c_i$ is a $k$-combination of monomials in the $z^{(c)}_j$, so $|c_i|_v\le1$. If
$|p|_v>1$, then $|c_ip^{n-i}|_v<|p|_v^n$ for all $i$, contradicting the ultrametric inequality. (2) By (1), $\mathcal X^\an$ lies in the
compact polydisc $\{|g_\ell|\le1\}$ for generators $g_\ell$ of $S$ (Lemma~\ref{lem:polydisc-compact}); multigradedness and normalization
are closed conditions on finitely many or countably many evaluations, so $\mathcal X^\an$ is compact; it is Hausdorff as a subspace of
$(\Spec S)^\an$. $\star$ is a multiplicative seminorm since $S$ is a domain, multigraded since a nonzero element has a nonzero
component, and normalized by (M1). If $k$ is countable, $S$ is countable and $(\Spec S)^\an$ is second countable.
\end{proof}

\begin{lemma}[Charts]\label{lem:gl-charts}
\lean{EnvelopeProperty.MultiModel.chartAlg, EnvelopeProperty.MultiModel.mem_chartAlg_iff, EnvelopeProperty.MultiModel.fg_chartAlg, EnvelopeProperty.MultiModel.isFractionRing_chartAlg, EnvelopeProperty.MultiModel.isIntegrallyClosed_chartAlg, EnvelopeProperty.MultiModel.chartDomain, EnvelopeProperty.MultiModel.chartHomeomorph, EnvelopeProperty.MultiModel.chartHomeomorph_triv, EnvelopeProperty.MultiModel.exists_chartIndex, EnvelopeProperty.MultiModel.iUnion_chartUnitDomain}
\leanok
\uses{def:gl-points, lem:gl-multigraded-extension, lem:bj-chart, lem:wu-proj-chart, lem:pg-chart-normal, lem:principal-open}
\source{fhx-envelope-20261028:sec:normal-local-regularization}
Let $j=(j_1,\dots,j_r)$ with $f^{(c)}_{j_c}\ne0$ for all $c$, $f_j=\prod_cf^{(c)}_{j_c}$, $z_j=\prod_cz^{(c)}_{j_c}$ and, for $a\in\N^r$,
$z_j^a=\prod_c(z^{(c)}_{j_c})^{a_c}$.
\begin{enumerate}
\item $A_j\coloneqq\{g/\tilde z_j^a:g\in S_a,\ a\in\N^r\}\subseteq K$ (where $z_j^a=\tilde z_j^aT^a$) is a finitely generated $k$-algebra with
fraction field $K$, generated by the $g_\ell/\tilde z_j^{d_\ell}$ for homogeneous generators $g_\ell\in S_{d_\ell}$ of $S$; it contains the
$f^{(c)}_i/f^{(c)}_{j_c}$, and it is integrally closed if $S$ is.
\item $U_j=\{|z_j|\ne0\}\subseteq\mathcal X^\an$ is open, and $|g/\tilde z_j^a|_{\iota_j(v)}=|gT^a|_v/|z_j^a|_v$ defines a homeomorphism
$\iota_j\colon U_j\to(\Spec A_j)^\an$ with $\iota_j(\star)=\vtriv$; the inverse sends $y$ to the point of
Lemma~\ref{lem:gl-multigraded-extension} with $N(gT^a)=|g/\tilde z_j^a|_y\prod_c\mu_c(y)^{-a_c}$,
$\mu_c(y)=\max_i|f^{(c)}_i/f^{(c)}_{j_c}|_y\ge1$.
\item The sets $\{v:|z^{(c)}_{j_c}|_v=1\ \forall c\}\subseteq U_j$ cover $\mathcal X^\an$.
\end{enumerate}
\end{lemma}

\begin{proof}
(1) A degree-zero fraction $p/z_j^a$, $p\in S_a T^a$, expands $p$ in monomials of the generators $g_\ell T^{d_\ell}$; a monomial of degree $a$
gives a product of the $g_\ell/\tilde z_j^{d_\ell}$. The fraction field is $K$ by (M3), since $f/g=(f/\tilde z_j^a)/(g/\tilde z_j^a)$.
$f^{(c)}_i/f^{(c)}_{j_c}=z^{(c)}_i/z^{(c)}_{j_c}$. Normality: if $w\in K$ is integral over $A_j$, then $w$ is integral over $S[z_j^{-1}]$,
which is integrally closed when $S$ is, so $wT^0\in S[z_j^{-1}]$ is a degree-zero element of $S[z_j^{-1}]$, i.e.\ $w\in A_j$ (the
argument of Lemma~\ref{lem:pg-chart-normal}). (2) $\iota_j(v)$ is a multiplicative seminorm, trivial on $k$, continuous in $v$. For the inverse,
$\mu_c(y)\ge1$ because $f^{(c)}_{j_c}/f^{(c)}_{j_c}=1$; $N$ satisfies the hypotheses of Lemma~\ref{lem:gl-multigraded-extension} ($N(z^{(c)}_i)=
|f^{(c)}_i/f^{(c)}_{j_c}|_y/\mu_c(y)$), and the two constructions are mutually inverse and continuous, as in Lemma~\ref{lem:bj-chart}(2).
$\iota_j(\star)$ is $1$ on nonzero elements. (3) For each $c$ choose $j_c$ with $|z^{(c)}_{j_c}|_v=1$.
\end{proof}

\begin{definition}[Sections]\label{def:gl-sections}
\lean{EnvelopeProperty.MultiModel.addDeg, EnvelopeProperty.MultiModel.IsAdmissible, EnvelopeProperty.MultiModel.sections, EnvelopeProperty.MultiModel.secAbsAt, EnvelopeProperty.MultiModel.secAbs, EnvelopeProperty.MultiModel.logSecAbs}
\leanok
\uses{def:gl-model, def:gl-points}
\source{fhx-envelope-20261028:eq:FS-function, bj22-trivval:defi:FS}
For $a\in\mathbf Z^r$, $R_a=R_a(\mathcal X)$ is the set of $f\in K$ for which there is $b\in\N^r$ with $a+b\in\N^r$ and
$f\tilde z^\beta\in S_{a+b}$ for all $\beta\in\mathrm{Mon}(b)$ (such $b$ is \emph{admissible} for $f$). For $f\in R_a$ and
$v\in\mathcal X^\an$ put $|f|_v\coloneqq\max_{\beta\in\mathrm{Mon}(b)}|f\tilde z^\beta T^{a+b}|_v$ and $\log|f|(v)=\log|f|_v$.
\end{definition}

\begin{lemma}[Properties of sections]\label{lem:gl-sections}
\lean{EnvelopeProperty.MultiModel.secAbs_eq, EnvelopeProperty.MultiModel.piece_le_sections, EnvelopeProperty.MultiModel.secAbs_of_mem_piece, EnvelopeProperty.MultiModel.mul_mem_sections, EnvelopeProperty.MultiModel.secAbs_mul, EnvelopeProperty.MultiModel.secAbs_add_le, EnvelopeProperty.MultiModel.secAbs_le_one, EnvelopeProperty.MultiModel.secAbs_triv, EnvelopeProperty.MultiModel.continuous_secAbs, EnvelopeProperty.MultiModel.secAbs_ne_zero, EnvelopeProperty.MultiModel.secAbs_eq_chartMap}
\leanok
\uses{def:gl-sections, lem:gl-points-basic, lem:gl-charts}
\begin{enumerate}
\item $|f|_v$ does not depend on the admissible $b$; for $a\in\N^r$, $S_a\subseteq R_a$ and $|f|_v=|fT^a|_v$.
\item $R_aR_{a'}\subseteq R_{a+a'}$, $|ff'|_v=|f|_v|f'|_v$, $|f+f'|_v\le\max(|f|_v,|f'|_v)$ for $f,f'\in R_a$, and $|\lambda f|_v=|f|_v$ for
$\lambda\in k^\times$.
\item $|f|_v\le1$, $|f|_\star=1$ for $f\ne0$, $v\mapsto|f|_v$ is continuous, and $|f|_v\ne0$ for $f\ne0$ and $v\in\mathcal X^{\mathrm{val}}$.
\item For a chart $j$ (Lemma~\ref{lem:gl-charts}), $f/\tilde z_j^a\in A_j$ (with $\tilde z_j^a\coloneqq\prod_c(f^{(c)}_{j_c})^{a_c}$) and
$|f|_v=|f/\tilde z_j^a|_{\iota_j(v)}\,|z_j^a|_v$ for $v\in U_j$, the right side read with $|z_j^a|_v=\prod_c|z^{(c)}_{j_c}|_v^{a_c}$.
\end{enumerate}
\end{lemma}

\begin{proof}
(1) If $b$ is admissible, so is every $b'\ge b$: $f\tilde z^{\beta+\gamma}=(f\tilde z^\beta)\tilde z^\gamma$. Every $\beta'\in\mathrm{Mon}(b')$ is $\beta+\gamma$
with $\beta\in\mathrm{Mon}(b)$, $\gamma\in\mathrm{Mon}(b'-b)$, so by multiplicativity
$\max_{\beta'}|f\tilde z^{\beta'}T^{a+b'}|=\max_\beta|f\tilde z^\beta T^{a+b}|\cdot\max_\gamma|z^\gamma|$, and $\max_\gamma|z^\gamma|=\prod_c(\max_j|z^{(c)}_j|)^{b'_c-b_c}=1$.
Two admissible vectors are compared through their maximum. For $f\in S_a$ take $b=0$. (2) If $b,b'$ are admissible for $f,f'$ then
$b+b'$ is admissible for $ff'$, and $\mathrm{Mon}(b+b')=\mathrm{Mon}(b)+\mathrm{Mon}(b')$ gives multiplicativity; the ultrametric inequality is
termwise. (3) $|f|\le1$ by Lemma~\ref{lem:gl-points-basic}(1); at $\star$ every $f\tilde z^\beta T^{a+b}$ with $\tilde z^\beta\ne0$ has norm $1$, and some
$\tilde z^\beta\ne0$ by (M1); a finite maximum of continuous functions is continuous; at a point with generic support every nonzero
element of $S$ has nonzero norm. (4) $f\tilde z_j^b\in S_{a+b}$ for $\beta=bj$ (the monomial $z_j^b$), so $f/\tilde z_j^a=(f\tilde z_j^b)/\tilde z_j^{a+b}\in A_j$;
for $v\in U_j$ normalize with the chart formula of Lemma~\ref{lem:gl-charts}(2).
\end{proof}

\begin{definition}[Classes and cones]\label{def:gl-cones}
\lean{EnvelopeProperty.MultiModel.IsBasePointFree, EnvelopeProperty.MultiModel.semiampleCone, EnvelopeProperty.MultiModel.nefCone, EnvelopeProperty.ratClosure}
\leanok
\uses{def:gl-sections}
The \emph{model-semiample cone} $C(\mathcal X)\subseteq\Q^r$ is the set of $a$ for which there are $m\ge1$ with $ma\in\mathbf Z^r$ and
finitely many $f_t\in R_{ma}$ without common zero on $\mathcal X^\an$. The \emph{model-nef cone} $N(\mathcal X)\subseteq\R^r$ is the
closure of $C(\mathcal X)$. A class $\theta$ is \emph{model-ample} if it is positive (all coordinates $>0$).
\end{definition}

\begin{lemma}[The cones]\label{lem:gl-cones}
\lean{EnvelopeProperty.MultiModel.nonneg_mem_semiampleCone, EnvelopeProperty.MultiModel.isClosed_nefCone, EnvelopeProperty.MultiModel.IsBasePointFree.sup_secAbs_eq_one, EnvelopeProperty.Berk.span_eq_top_of_forall_exists_val_ne_zero, EnvelopeProperty.mem_of_lt_of_mem_ratClosure, EnvelopeProperty.exists_mem_ge_tendsto_of_mem_ratClosure}
\leanok
\uses{def:gl-cones, lem:gl-sections, lem:gl-charts}
\begin{enumerate}
\item $\Q^r_{\ge0}\subseteq C(\mathcal X)$, $C+C\subseteq C$, $\Q_{>0}C\subseteq C$ and $C+\Q^r_{\ge0}\subseteq C$; $N(\mathcal X)$ is a closed convex
cone containing $\R^r_{\ge0}$.
\item If $f_t\in R_{ma}$ have no common zero, then $\max_t|f_t|_v=1$ for every $v$.
\item If $\theta\in N(\mathcal X)$ and $(a_\alpha)_{\alpha\in I}$ is a net in $\Q^r$ converging to $\theta$, there is a net $(a'_{\alpha,n})$ indexed by $I\times\N$
(product order) with $a'_{\alpha,n}\in C(\mathcal X)$, $a'_{\alpha,n}\ge a_\alpha$ and $a'_{\alpha,n}\to\theta$. (Re-indexing is needed: for a constant one-point net
$a=\theta$ with $\theta\in N(\mathcal X)\cap\Q^r$ outside $C(\mathcal X)$, e.g.\ a degree-$0$ non-torsion class on an elliptic curve, no net in $C$ indexed by $I$
converges to $\theta$ while dominating $a$.)
\end{enumerate}
\end{lemma}

\begin{proof}
(1) Monomials $\tilde z^\beta$, $\beta\in\mathrm{Mon}(ma)$, have no common zero (some $|z^{(c)}_j|=1$); products and powers of families without common
zero have none; multiply by monomials. (2) Let $v\in U_j$ with $|z^{(c)}_{j_c}|_v=1$ (Lemma~\ref{lem:gl-charts}(3)). The elements
$g_t=f_t/\tilde z_j^{ma}\in A_j$ have no common zero on $(\Spec A_j)^\an$ (Lemma~\ref{lem:gl-sections}(4)), so they generate the unit ideal (a
prime containing them would give the trivial seminorm of $A_j/\mathfrak p$ as a common zero): $1=\sum c_tg_t$ with $c_t\in A_j$. Every element of
$A_j$ has norm $\le1$ at $\iota_j(v)$ (it is $p/z_j^a$ with $|p|_v\le1$, $|z_j|_v=1$), so $1\le\max_t|g_t|_{\iota_j(v)}=\max_t|f_t|_v\le1$. (3) Choose rationals $q_n\in\Q^r$ with $q_n>\theta$ coordinatewise and $q_n\to\theta$. Every rational class $q>\theta$ lies in $C$: some $c\in C$ is so close to
$\theta$ that $q-c\ge0$, and $C+\Q^r_{\ge0}\subseteq C$. Put $a'_{\alpha,n}=a_\alpha\vee q_n$ (coordinatewise maximum). Then $a'_{\alpha,n}>\theta$, so $a'_{\alpha,n}\in C$;
$a'_{\alpha,n}\ge a_\alpha$; and $a'_{\alpha,n}\to\theta\vee\theta=\theta$ along $I\times\N$.
\end{proof}

\begin{definition}[Fubini--Study functions]\label{def:gl-fs}
\lean{EnvelopeProperty.MultiModel.FSData, EnvelopeProperty.MultiModel.FSData.eval, EnvelopeProperty.MultiModel.IsFS, EnvelopeProperty.MultiModel.IsFinFS, EnvelopeProperty.MultiModel.IsHomFS}
\leanok
\uses{def:gl-sections}
\source{fhx-envelope-20261028:eq:FS-function, bj22-trivval:defi:FS}
For $a\in\Q^r$, $\mathcal H(\mathcal X,a)$ is the set of functions $m^{-1}\max_t\{\log|f_t|+\lambda_t\}$ with $m\in\mathbf Z_{>0}$,
$ma\in\mathbf Z^r$, finitely many $f_t\in R_{ma}$ not all zero, and $\lambda_t\in\R$. $\mathcal H_{\mathrm{fin}}(\mathcal X,a)$ requires moreover
$\lambda_t\in\Q$ and no common zero on $\mathcal X^\an$, and $\mathcal H_{\mathrm{hom}}(\mathcal X,a)$ requires all $\lambda_t=0$.
\end{definition}

\begin{definition}[Fubini--Study systems with a semiample cone]\label{def:gl-fs-system}
\lean{EnvelopeProperty.FSSystemC, EnvelopeProperty.FSSystemC.coneClosure, EnvelopeProperty.FSSystemC.IsPSH, EnvelopeProperty.FSSystem.toFSSystemC, EnvelopeProperty.FSSystem.coneClosure_toFSSystemC}
\leanok
\uses{def:pg-fs-system}
\source{bj22-trivval:defi:FS, bj22-trivval:prop:FS}
Let $K$ be a compact space, $\star\in K$, and $C\subseteq\Q^r$ a convex cone with $\Q^r_{\ge0}\subseteq C$ and $C+\Q^r_{\ge0}\subseteq C$. A
\emph{Fubini--Study system with semiample cone $C$} assigns to every $a\in\Q^r$ sets $\mathcal H_{\mathrm{fin}}(a)\subseteq\mathcal H(a)$ of
continuous functions, $\mathcal H_{\mathrm{fin}}(a)$ real valued, satisfying (FS1), (FS2), (FS4) of Definition~\ref{def:pg-fs-system} for all
$a,b\in\Q^r$, and (FS3), (FS5) for $a\in C$ (constants, maxima with constants, approximation by $\mathcal H_{\mathrm{fin}}$); for every $a$,
$\mathcal H(a)$ is stable under finite maxima and addition of real constants. $\theta$-psh functions, $\theta\in\R^r$, are defined as in
Definition~\ref{def:pg-psh}. The systems of Definition~\ref{def:pg-fs-system} are those with $C=\Q^r_{\ge0}$ and $\mathcal H(a)=\varnothing$ unless
$a\ge0$.
\end{definition}

\begin{theorem}[Psh theory of a system]\label{thm:gl-fs-system-theory}
\lean{EnvelopeProperty.FSSystemC.IsFS.isPSH, EnvelopeProperty.FSSystemC.IsPSH.le_apply_triv, EnvelopeProperty.FSSystemC.IsPSH.max, EnvelopeProperty.FSSystemC.IsPSH.add, EnvelopeProperty.FSSystemC.IsPSH.smul_rat, EnvelopeProperty.FSSystemC.IsPSH.mono_class, EnvelopeProperty.FSSystemC.isPSH_of_antitone, EnvelopeProperty.FSSystemC.IsPSH.smul, EnvelopeProperty.FSSystemC.isPSH_const, EnvelopeProperty.FSSystemC.IsPSH.exists_isFinFS_net, EnvelopeProperty.FSSystemC.isPSH_iff_exists_isFinFS, EnvelopeProperty.FSSystemC.IsPSH.exists_isFinFS_net_of_pos, EnvelopeProperty.FSSystemC.IsPSH.exists_isFinFS_seq_gt}
\leanok
\uses{def:gl-fs-system, prop:pg-maxprin, thm:pg-psh2, thm:pg-psh-basic, thm:pg-regularization, cor:pg-countable-regularization, lem:bj-psh-diagonal, rmk:bj-psh-diagonal-finite}
\source{bj22-trivval:thm:psh12, bj22-trivval:thm:PSH, bj22-trivval:thm:pshample, bj22-trivval:prop:pshconvex}
Let $\overline C\subseteq\R^r$ be the closure of $C$.
\begin{enumerate}
\item For every $\theta\in\R^r$: $\mathcal H(a)\subseteq\PSH(a)$; every $u\in\PSH(\theta)$ satisfies $u(\star)\in\R$, $u\le u(\star)$; $\PSH(\theta)$ is
stable under finite maxima and real translates; $\PSH(\theta)+\PSH(\theta')\subseteq\PSH(\theta+\theta')$; $c\PSH(\theta)\subseteq\PSH(c\theta)$
for $c\in\Q_{>0}$; $\PSH(\theta)\subseteq\PSH(\theta')$ for $\theta\le\theta'$.
\item For $\theta\in\overline C$: PSH2 (Theorem~\ref{thm:pg-psh2}(3)) holds for nets $u_i\in\PSH(\theta_i)$, $\theta_i\to\theta$; real scaling
$t\PSH(\theta)=\PSH(t\theta)$; $\PSH(\theta)$ contains the real constants; the regularization statements of
Theorem~\ref{thm:pg-regularization} and Corollary~\ref{cor:pg-countable-regularization} hold, part~(2) for rational positive $\theta$.
\end{enumerate}
\end{theorem}

\begin{proof}
(1) These are the proofs of Proposition~\ref{prop:pg-maxprin} (without the conclusion $\theta\ge0$, which used $\mathcal H(a)=\varnothing$ for
$a\not\ge0$) and of Theorem~\ref{thm:pg-psh-basic}(1),(2),(3): maxima and sums use a common index set and the classes $a_i\vee a'_j$,
$a_i+a'_j$, and (FS2), (FS4); monotonicity uses $\mathcal H(a)\subseteq\mathcal H(a+(\theta'-\theta)_{\text{rounded up}})$ along the net. No constants
and no diagonal lemma are needed. (2) Let $u_i\in\PSH(\theta_i)$ decrease, $\theta_i\to\theta\in\overline C$. For $n\in\N$ put $\tau_{i,n}=\theta_i\vee\theta+(n+1)^{-1}\mathbf 1$; then
$\tau_{i,n}>\theta$, $\tau_{i,n}\to\theta$ along $I\times\N$, and $u_i\in\PSH(\tau_{i,n})$ by monotonicity (1). The classes of the approximating nets of
$u_i$ in $\PSH(\tau_{i,n})$ converge to $\tau_{i,n}>\theta$, so they are eventually $>\theta$ and hence in $C$ (as in the proof of Lemma~\ref{lem:gl-cones}(3)).
(Replacing the classes $a_{i,\alpha}$ of the original approximating nets does not work: they converge to $\theta_i$, which need not lie in $\overline C$.) The families $\mathcal H(a)$, $a\in C$, consist of continuous functions avoiding $+\infty$ and are stable
under addition of constants and maxima with constants (FS3), so the diagonal lemma in the form of Remark~\ref{rmk:bj-psh-diagonal-finite}
(\texttt{IsDecrNetLimit.of\_antitone}) applies and gives PSH2. The other statements are proved as in Theorems~\ref{thm:pg-psh-basic},
\ref{thm:pg-regularization} and Corollary~\ref{cor:pg-countable-regularization}, after the same replacement of classes.
\end{proof}

\begin{lemma}[The system of a model]\label{lem:gl-fs-system}
\lean{EnvelopeProperty.MultiModel.fsSystem, EnvelopeProperty.MultiModel.coneClosure_fsSystem, EnvelopeProperty.MultiModel.IsFinFS.mem_semiampleCone, EnvelopeProperty.MultiModel.isFS_const, EnvelopeProperty.MultiModel.IsFS.exists_antitone_isFinFS, EnvelopeProperty.MultiModel.isPSH_logSecAbs}
\leanok
\uses{def:gl-fs, def:gl-fs-system, lem:gl-sections, lem:gl-cones, lem:pg-subvariety-fs-system}
$(\mathcal X^\an,\star)$ with $\mathcal H(\mathcal X,a)\supseteq\mathcal H_{\mathrm{fin}}(\mathcal X,a)$ is a Fubini--Study system with semiample
cone $C(\mathcal X)$.
\end{lemma}

\begin{proof}
As in Lemma~\ref{lem:pg-subvariety-fs-system}, with sections in place of forms. (FS1): $|f|\le1=|f|_\star$ for $0\ne f$ (Lemma~\ref{lem:gl-sections}(3)),
so $\varphi\le\varphi(\star)=m^{-1}\max\{\lambda_t:f_t\ne0\}$. (FS2): for $b-a\ge0$, pass to a common $m$ with $m(b-a)\in\N^r$ and replace the family by
$f_t\tilde z^\gamma$, $\gamma\in\mathrm{Mon}(m(b-a))$; since $\max_\gamma|\tilde z^\gamma|=1$ the function is unchanged, nonzero and base point freeness
are preserved. (FS3) for $a\in C$: a constant $c$ is $m^{-1}\max_t\{\log|f_t|+mc\}$ for a base point free family (Lemma~\ref{lem:gl-cones}(2));
maxima with constants are unions of families. (FS4): products $f_t^{m'}f'^{\,m}_u$ and powers. (FS5) for $a\in C$: rational weights
$\lambda^{(n)}_t\downarrow\lambda_t$ and truncation $\max\{\cdot,-n\}$ with the constant's base point free family.
\end{proof}

\begin{definition}[Psh functions and the envelope property]\label{def:gl-psh}
\lean{EnvelopeProperty.MultiModel.IsPSH, EnvelopeProperty.MultiModel.envelope, EnvelopeProperty.MultiModel.HasEnvelopeProperty, EnvelopeProperty.MultiModel.IsPSHHom}
\leanok
\uses{lem:gl-fs-system, thm:gl-fs-system-theory, def:pn-ep}
\source{fhx-envelope-20261028:subsec:prelim-psh, bj22-trivval:defi:psh, bj22-trivval:sec:envprop}
For $\theta\in\R^r$, $\PSH(\mathcal X,\theta)$ is the set of $\theta$-psh functions of the system of Lemma~\ref{lem:gl-fs-system}:
usc functions $u\not\equiv-\infty$ that are limits of decreasing nets $\varphi_\alpha\in\mathcal H(\mathcal X,a_\alpha)$, $a_\alpha\to\theta$.
$\PSH_{\mathrm{hom}}(\mathcal X,\theta)$ is the subset of homogeneous functions (Definition~\ref{def:gl-scaling}). $\theta$ has the
\emph{envelope property}, $\mathrm{EP}_{\mathcal X}(\theta)$, if $(\sup_iu_i)^\star\in\PSH(\mathcal X,\theta)$ for every nonempty family
$(u_i)\subseteq\PSH(\mathcal X,\theta)$ bounded above. For $\phi\colon\mathcal X^\an\to\R$ bounded, $\mathrm P_\theta(\phi)=\sup\{\tau\in
\PSH(\mathcal X,\theta):\tau\le\phi\}$.
\end{definition}

\begin{lemma}[Elementary rules for the envelope property]\label{lem:gl-ep-basic}
\lean{EnvelopeProperty.MultiModel.HasEnvelopeProperty.of_tendsto, EnvelopeProperty.MultiModel.hasEnvelopeProperty_of_forall_rat_gt, EnvelopeProperty.MultiModel.hasEnvelopeProperty_smul_iff, EnvelopeProperty.MultiModel.hasEnvelopeProperty_of_forall_not_isPSH, EnvelopeProperty.FSSystemC.HasEnvelopeProperty.of_tendsto}
\leanok
\uses{def:gl-psh, thm:gl-fs-system-theory, lem:bj-envppty}
\source{bj22-trivval:lem:envppty, fhx-envelope-20261028:thm:normal-ample-EP}
\begin{enumerate}
\item If $\theta\in N(\mathcal X)$ is the limit of classes $\theta_j\ge\theta$ with $\mathrm{EP}(\theta_j)$, then $\mathrm{EP}(\theta)$.
\item For $\theta\in N(\mathcal X)$ and real $t>0$, $\mathrm{EP}(\theta)\iff\mathrm{EP}(t\theta)$.
\item If $\PSH(\mathcal X,\theta)=\varnothing$, $\mathrm{EP}(\theta)$ holds.
\end{enumerate}
\end{lemma}

\begin{proof}
(1) As Lemma~\ref{lem:bj-envppty}: $u=(\sup u_i)^\star\in\PSH(\theta_j)$ for all $j$ by monotonicity, and PSH2 at $\theta\in N(\mathcal X)$
(Theorem~\ref{thm:gl-fs-system-theory}(2)) for the constant sequence. (2) $t\PSH(\theta)=\PSH(t\theta)$ and $(\sup tu_i)^\star=t(\sup u_i)^\star$. (3) There
is no nonempty family.
\end{proof}

\begin{remark}[The model and $N^1(X)_\R$]\label{rmk:gl-model-comparison}
\uses{def:gl-psh, def:gl-cones, rmk:pg-subvariety-psh-comparison, rmk:pg-blowup-psh-comparison}
Let $X$ be an integral projective variety. Choose globally generated line bundles $L_1,\dots,L_r$ whose classes span $N^1(X)_\Q$, one of them
very ample, and generating sections; their product map embeds $X$ as $V(I)\subseteq\prod_c\PP^{M_c}$, giving a quotient model
(Remark~\ref{rmk:gl-model-examples}(1)) with $c_1\colon\R^r\to N^1(X)_\R$, $\theta\mapsto\sum_c\theta_cc_1(L_c)$, surjective. For this model the
following are not formalized: (a)
$\mathcal X^\an=X^\an$ (as Definition~\ref{def:pn-an}); (b) $R_a=H^0(X,\sum_ca_cL_c)$ (the ideal transform of $S$ with respect to the irrelevant
ideal), so $\mathcal H(\mathcal X,a)=\mathcal H^{\mathrm{gf}}_\R(\sum_ca_cL_c)$; (c) $C(\mathcal X)$ consists of the rational vectors whose class is
semiample, hence $N(\mathcal X)=c_1^{-1}(\mathrm{Nef}(X))$ (ample classes are dense in the nef cone, and $\theta+\delta\mathbf 1$ has ample class for
$\theta\in c_1^{-1}\mathrm{Nef}$, $\delta>0$); (d) $\PSH(\mathcal X,\theta)=\PSH(X,c_1(\theta))$ of \cite[Definition~4.1]{BJ22}: approximants of
\cite{BJ22} with classes $c_1(L_i)\to c_1(\theta)$ lie in $\mathcal H(\mathcal X,a)$ for every rational $a\ge\theta+\delta\mathbf1$ and $i\gg0$, since
$c_1(a)-c_1(L_i)$ is then ample (openness of the ample cone, $c_1(\mathbf 1)$ ample); (e) a modification or a finite cover of $X$ is a model in the
same coordinates (Remark~\ref{rmk:gl-model-examples}(2),(3)), so pulled-back classes need no new coordinates (this is
\cite[Lemma~2.6]{FHX}, \texttt{lem:finite-class-promotion}, in the model). Internally, ampleness is used only through the sufficient condition
``positive'', and nefness through $C(\mathcal X)$; the shape of the nef cone (simplicial or not, polyhedral or not) plays no role. Every
$\theta\in N^1(X)_\R$ is $c_1(\theta)$ for some $\theta\in\R^r$, possibly with negative coordinates, so the formal statements for all
$\theta\in\R^r$ (resp.\ $\theta\in N(\mathcal X)$) are the statements of \cite{FHX} for all real (resp.\ all nef) classes.
The statements of \cite{FHX} for $X$ itself and every $\theta\in N^1(X)_\R$ are formalized in Section~\ref{sec:source-main}
through a comparison that chooses no line bundles spanning $N^1(X)_\Q$: every class is $\pi(\theta_m)$ for a presentation
of $X$ (Lemma~\ref{lem:pm-exists-presentation}), and the psh functions and the envelope property of $X$ are those of the
presentation model (Theorem~\ref{thm:var-psh-comparison}, Corollary~\ref{cor:var-ep-comparison}); this gives
Theorem~\ref{thm:var-main} and Corollary~\ref{cor:var-characterization}, with ``ample'' in~(3) meaning membership in the
ample cone.
\end{remark}

\begin{definition}[Morphisms of models]\label{def:gl-maps}
\lean{EnvelopeProperty.MultiModel.FiniteHom, EnvelopeProperty.MultiModel.FiniteHom.anMap, EnvelopeProperty.MultiModel.FiniteHom.finiteDimensional, EnvelopeProperty.MultiModel.ScalarRestriction, EnvelopeProperty.MultiModel.ScalarRestriction.anMap, EnvelopeProperty.MultiModel.FiniteHom.ofNormalization, EnvelopeProperty.MultiModel.FiniteHom.ofNormalizationIn}
\leanok
\uses{def:gl-model, def:gl-points, def:gl-sections}
\source{fhx-envelope-20261028:prop:BJ-package, fhx-envelope-20261028:subsec:prelim-finite-classes}
A \emph{finite morphism} $h\colon\mathcal T\to\mathcal X$ of models of rank $r$ is an inclusion $K\subseteq E$ of finite degree $d$ together with
$S\subseteq S_T\subseteq E[T]$, $S_T$ integral over $S$, with the same coordinates; $h^\an(w)=w|_S$. A \emph{scalar restriction} along
$k_0\subseteq k$ ($k_0$ algebraically closed) is a model $\mathcal X_0=(K_0,S_0,z)$ over $k_0$ and $\mathcal X=(K,S,z)$ over $k$ with
$S=k\cdot S_0\cong k\otimes_{k_0}S_0$ inside $K[T]$; $r(v)=v|_{S_0}$.
\end{definition}

\begin{lemma}[Pullbacks]\label{lem:gl-pullback}
\lean{EnvelopeProperty.MultiModel.FiniteHom.continuous_anMap, EnvelopeProperty.MultiModel.FiniteHom.anMap_surjective, EnvelopeProperty.MultiModel.FiniteHom.finite_anMap_preimage, EnvelopeProperty.MultiModel.FiniteHom.eq_triv_of_anMap_eq_triv, EnvelopeProperty.MultiModel.FiniteHom.secAbs_anMap, EnvelopeProperty.MultiModel.FiniteHom.isFS_comp, EnvelopeProperty.MultiModel.FiniteHom.isPSH_comp, EnvelopeProperty.MultiModel.ScalarRestriction.continuous_anMap, EnvelopeProperty.MultiModel.ScalarRestriction.anMap_triv, EnvelopeProperty.MultiModel.ScalarRestriction.secAbs_anMap, EnvelopeProperty.MultiModel.ScalarRestriction.isFS_comp, EnvelopeProperty.MultiModel.ScalarRestriction.isPSH_comp}
\leanok
\uses{def:gl-maps, lem:gl-sections, def:gl-psh}
\source{fhx-envelope-20261028:prop:BJ-package, bj22-trivval:thm:PSH}
For a finite morphism $h$: $h^\an$ is continuous, surjective and has finite fibres; $R_a(\mathcal X)\subseteq R_a(\mathcal T)$ with the same values,
so $\varphi\circ h^\an\in\mathcal H(\mathcal T,a)$ for $\varphi\in\mathcal H(\mathcal X,a)$ and $u\circ h^\an\in\PSH(\mathcal T,\theta)$ for
$u\in\PSH(\mathcal X,\theta)$; the only point of $\mathcal T^\an$ over $\star$ is $\star$. For a scalar restriction, $r$ is continuous,
$r(\star)=\star$, values of sections are preserved ($R_a(\mathcal X_0)\subseteq R_a(\mathcal X)$), and FS and psh functions pull back; surjectivity,
finiteness of fibres and uniqueness over $\star$ fail in general and are not needed.
\end{lemma}

\begin{proof}
$h^\an(w)$ is multigraded and normalized, as $S\subseteq S_T$ is graded with the same coordinates. Surjectivity and finiteness of fibres: on a chart
$j$, $A_j\subseteq B_j$ (the chart algebra of $\mathcal T$) is finite, and points of $(\Spec B_j)^\an$ over a point of $(\Spec A_j)^\an$ exist and are finite
in number (integral extension of seminorms; fibres $\leftrightarrow$ maximal ideals of $\mathcal H(x)\otimes_{A_j}B_j$, Lemma~\ref{lem:finite-flat-norm}
and \texttt{Berkovich/Finite}); charts correspond since coordinates agree. Admissible $b$ for $f\in R_a(\mathcal X)$ stay admissible in $S_T$, and
$|f|_w=|f|_{h(w)}$. Pullbacks of decreasing nets are decreasing nets with the same classes, and $u\circ h^\an\not\equiv-\infty$ since
$u(\star)>-\infty$ and $h(\star)=\star$. A point over $\star$ is $1$ on $S\setminus0$; an element of $S_T\setminus0$ satisfies an integral equation over $S$
whose constant term may be taken nonzero, which forces norm $1$. For a scalar restriction only the value formula
$|f|_{r(v)}=|f|_v$ and $r(\star)=\star$ are used. The other assertions fail: for $k_0=\overline{\Q}\subseteq k=\overline{\Q(t)}$ and
$\mathcal X=\PP^1$, the divisorial point $e^{-\mathrm{ord}_{x=t}}$ is not $\star$ but restricts to $\star$ on $k_0[x]$, and so does the
point at $x=s$ for every $s$ transcendental over $k_0$.
\end{proof}

\begin{lemma}[Change of model]\label{lem:gl-model-change}
\lean{EnvelopeProperty.MultiModel.IsModelChange, EnvelopeProperty.MultiModel.IsModelChange.restrictHomeomorph, EnvelopeProperty.MultiModel.IsModelChange.hasSections, EnvelopeProperty.MultiModel.IsModelChange.sections_eq, EnvelopeProperty.MultiModel.IsModelChange.secAbs_restrictAn, EnvelopeProperty.MultiModel.IsModelChange.isPSH_comp_iff, EnvelopeProperty.MultiModel.IsModelChange.hasEnvelopeProperty_iff, EnvelopeProperty.MultiModel.IsModelChange.semiampleCone_eq, EnvelopeProperty.MultiModel.IsModelChange.nefCone_eq, EnvelopeProperty.MultiModel.isPSH_normalization_iff, EnvelopeProperty.MultiModel.hasEnvelopeProperty_normalization_iff}
\leanok
\uses{def:gl-model, lem:gl-charts, lem:gl-sections, def:gl-psh, lem:gl-points-basic}
\source{fhx-envelope-20261028:prop:BJ-package}
Let $\mathcal X=(K,S,z)$ and $\mathcal X'=(K,S',z')$ be models of rank $r$ with the same $K$ such that every coordinate of either model is
integral over the coordinates of the other in the same degree, and $S\subseteq S'\subseteq S^\nu$. Suppose the chart algebras of $\mathcal X$ are
integrally closed. Then restriction $v'\mapsto v'|_S$ is a homeomorphism $\mathcal X'^\an\to\mathcal X^\an$, $R_a(\mathcal X)=R_a(\mathcal X')$ with
equal values, and $\PSH(\mathcal X,\theta)=\PSH(\mathcal X',\theta)$, $C(\mathcal X)=C(\mathcal X')$.
\end{lemma}

\begin{proof}
Normalizations agree: $\max_j|z'^{(c)}_j|=1$ iff $\max_j|z^{(c)}_j|=1$, by the monic root estimate of Lemma~\ref{lem:gl-points-basic}(1) in both
directions. An element of $S^\nu$ of degree $a$ divided by $\tilde z_j^a$ is integral over the integrally closed $A_j$, hence lies in $A_j$; so for each
chart $f\tilde z_j^{c_j}\in S$ for some $c_j$. The homogeneous ideal $\{q\in S:fq\in S\}$ thus contains a power of every $z_j$; a prime containing all $z_j$
contains a whole block of coordinates, so its radical contains the product of the block ideals $(z^{(c)}_0,\dots,z^{(c)}_{M_c})$, and it contains $\mathrm{Mon}(b)$ for
$b\gg0$: $S^\nu_a\subseteq R_a(\mathcal X)$ (the same ``ideal argument'' is used repeatedly below). The two section spaces are defined
with the monomials $\tilde z^\beta$ and $\tilde z'^\gamma$ respectively; by mutual integrality, a power of each coordinate of one model lies in
the $S'$-ideal generated by the coordinates of the other model in the same block, so the ideal argument (with explicit exponents) gives
both inclusions, $R(\mathcal X')=R(\mathcal X)$, and the values agree (a double maximum over $\beta$ and $\gamma$, both points being
normalized for both coordinate systems). Then a point of $\mathcal X^\an$ extends to $S'$ by $|f|=\max_\beta|f\tilde z^\beta|$, uniquely and continuously; restriction is a
continuous bijection of compact Hausdorff spaces. The Fubini--Study families, hence psh functions and the semiample cones, coincide.
\end{proof}

\begin{lemma}[Veronese and Segre--Veronese models]\label{lem:gl-veronese}
\lean{EnvelopeProperty.MultiModel.veronese, EnvelopeProperty.MultiModel.veroneseHomeomorph, EnvelopeProperty.MultiModel.veroneseHomeomorph_triv, EnvelopeProperty.MultiModel.isNormal_veronese, EnvelopeProperty.MultiModel.mem_veronese_sections_iff, EnvelopeProperty.MultiModel.secAbs_veroneseAnMap, EnvelopeProperty.MultiModel.isFS_veronese_iff, EnvelopeProperty.MultiModel.isPSH_veronese_iff, EnvelopeProperty.MultiModel.hasEnvelopeProperty_veronese_iff, EnvelopeProperty.MultiModel.multiVeronese, EnvelopeProperty.MultiModel.multiVeroneseHomeomorph, EnvelopeProperty.MultiModel.isPSH_multiVeronese_iff, EnvelopeProperty.MultiModel.hasEnvelopeProperty_multiVeronese_iff, EnvelopeProperty.MultiModel.biVeronese}
\leanok
\uses{def:gl-model, lem:gl-multigraded-extension, def:gl-psh, thm:gl-fs-system-theory, lem:wu-veronese}
\source{fhx-envelope-20261028:thm:normal-ample-EP}
Let $a\in\N^r$ be positive. The \emph{Segre--Veronese model} $\mathcal X_{(a)}$ of rank $1$ has $S_{(a)}=\bigoplus_nS_{na}U^n\subseteq K[U]$ and
coordinates $\tilde z^\beta U$, $\beta\in\mathrm{Mon}(a)$. Restriction is a homeomorphism $\mathcal X^\an\to\mathcal X_{(a)}^\an$, preserving $\star$ and
values of sections ($R_n(\mathcal X_{(a)})=R_{na}(\mathcal X)$), and $\PSH(\mathcal X_{(a)},t)=\PSH(\mathcal X,ta)$ for every real $t$. $\mathcal X_{(a)}$ is
normal if $\mathcal X$ is. Hence $\mathrm{EP}_{\mathcal X}(ta)\iff\mathrm{EP}_{\mathcal X_{(a)}}(t)$.
\end{lemma}

\begin{proof}
$S_{(a)}$ is integral over the coordinates (homogeneous integral equations of degrees in $\N a$ have coefficients in $k[z]_{\N a}$, spanned by products of
monomials of degree $a$). A point of $\mathcal X^\an$ restricts to a point (normalization: $\max_\beta|z^\beta|=1$). Conversely, for $y\in\mathcal X_{(a)}^\an$
and $p\in S_b$ homogeneous, choose $n$ with $na\ge b$ and put $N(pT^b)=\max_{\gamma\in\mathrm{Mon}(na-b)}|p\tilde z^\gamma U^n|_y$; this is independent of $n$
(as in Lemma~\ref{lem:gl-sections}(1)) and satisfies the hypotheses of Lemma~\ref{lem:gl-multigraded-extension} (normalization:
$N(z^{(c)}_j)=\max_\gamma|z^{(c)}_j\tilde z^\gamma U^n|_y$ and $\max_j$ of it is $\max_{\beta\in\mathrm{Mon}(na)}|z^\beta|=1$), and gives the inverse. FS
functions of class $t\in\Q$ on $\mathcal X_{(a)}$ are those of class $ta$ on $\mathcal X$ (for data of class $q_0a$ on $\mathcal X$ with
multiplier $m$, the degree $mq_0$ on $\mathcal X_{(a)}$ need not be an integer: replace $m$ by $a_0m$ and the sections $f$ by $f^{a_0}$, where
$a_0$ is the denominator of $q_0$). Psh: an approximant of class $b_\alpha\to ta$ on $\mathcal X$ lies
in $\mathcal H(\mathcal X,\lambda_\alpha a)$ with $\lambda_\alpha=\max_cb_{\alpha,c}/a_c\to t$ (FS2), which is a class of $\mathcal X_{(a)}$. Normality: the argument of
Lemma~\ref{lem:wu-veronese} (second part), componentwise in $\N^r$.
\end{proof}

\begin{remark}[Bridges to the projective models already formalized]\label{rmk:gl-bridges}
\uses{rmk:gl-model-examples, lem:gl-model-change, lem:gl-veronese, lem:gl-bridge-points, lem:gl-bridge-psh, def:pg-subvariety-psh, def:pg-blowup-psh, def:pn-psh}
For $r=1$ and $S$ generated in degree $1$, a presentation $k[x_0,\dots,x_N]\to S$, $x_i\mapsto w_i$ (a basis of $S_1$), identifies $\mathcal X$ with the
quotient model of the kernel, whose points are those of $V(I)\subseteq\PP^N$ (Lemma~\ref{lem:gl-bridge-points}); for $r=2$ the same holds for
$\mathrm{BlowupCenter}$ models. The model sections $R_a$ contain the polynomial sections $S_a$ but can be larger (they are the sections of
$\mathcal O(a)$ on $\operatorname{Proj}S$), so the identification of psh functions (Lemma~\ref{lem:gl-bridge-psh}) needs an argument: $R_a=S_a$
when $S$ is integrally closed and generated in degree $1$, and otherwise integrality of $fT^a$ over $S$ and PSH2. By Lemma~\ref{lem:wu-veronese}
every rank-$1$ model has a Veronese $S_{(e)}$ generated in degree $1$, and Lemma~\ref{lem:gl-veronese} applies.
\end{remark}

\begin{lemma}[Bridges: points]\label{lem:gl-bridge-points}
\lean{EnvelopeProperty.MultiModel.DegOnePres.anHomeomorph, EnvelopeProperty.MultiModel.DegOnePres.anHomeomorph_val, EnvelopeProperty.MultiModel.DegOnePres.anHomeomorph_triv, EnvelopeProperty.MultiModel.DegOnePres.recoordHomeomorph, EnvelopeProperty.MultiModel.DegOnePres.isPSH_recoord_iff, EnvelopeProperty.MultiModel.DegOnePres.hasEnvelopeProperty_recoord_iff, EnvelopeProperty.MultiModel.biAnHomeomorph, EnvelopeProperty.MultiModel.blowupAnHomeomorph, EnvelopeProperty.MultiModel.blowupAnHomeomorph_triv, EnvelopeProperty.MultiModel.proj_blowupAnHomeomorph_val}
\leanok
\uses{rmk:gl-model-examples, lem:gl-model-change, lem:gl-points-basic, def:pg-subvariety-fs, def:pg-bigraded, def:pg-blowup}
\begin{enumerate}
\item Let $r=1$ and let $w_0,\dots,w_N\in S_1$ generate $S$, with kernel $I$ of $k[x_0,\dots,x_N]\to S$. Restriction along $x_i\mapsto w_iT$ is a
homeomorphism $\tau\colon V(I)^\an\to\mathcal X^\an$ with $\tau(\star)=\star$ and $|F(w)|_{\tau(u)}=|F|_u$; changing the coordinates of $\mathcal X$ to
$w$ does not change points, values of sections, psh functions or the envelope property.
\item Let $r=2$ and $S=k[z]$, with kernel $I'$ of $k[x,y]\to S$. Restriction is a homeomorphism of $\mathcal X^\an$ onto $V(I')^\an\subseteq
(\PP^M\times\PP^{M'})^\an$; if $I'$ is the ideal of a blowup $B$ (Definition~\ref{def:pg-blowup}), it is a homeomorphism $\vartheta\colon
\mathcal X^\an\to B^\an$ with $\vartheta(\star)=\star$, compatible with $\mu$.
\end{enumerate}
\end{lemma}

\begin{proof}
A point of $V(I)^\an$ is a multiplicative seminorm on $k[x]/I$ normalized by $\max_i|x_i|=1$; it is a seminorm on $S$, and normalization for the
coordinates $z$ and for $w$ agree by the monic root estimate (Lemma~\ref{lem:gl-points-basic}(1)), as each set of coordinates is integral over the
other in degree $1$. Continuity in both directions holds on generators. The second part is the same with bihomogeneous elements.
\end{proof}

\begin{lemma}[Bridges: sections and psh functions]\label{lem:gl-bridge-psh}
\lean{EnvelopeProperty.MultiModel.isPSH_comp_blowupAnHomeomorph_iff, EnvelopeProperty.MultiModel.isPSH_iff_isPSH_comp_blowupAnHomeomorph_symm, EnvelopeProperty.MultiModel.isFS_comp_blowupAnHomeomorph, EnvelopeProperty.MultiModel.IsFS.exists_isFS_blowup, EnvelopeProperty.MultiModel.IsPSH.blowup_comp_symm, EnvelopeProperty.MultiModel.hasEnvelopeProperty_iff_blowup, EnvelopeProperty.MultiModel.IsAdmissible.pow_of_eq_coordAlg, EnvelopeProperty.ProjSubvariety.BlowupCenter.multiModel, EnvelopeProperty.ProjSubvariety.BlowupCenter.multiModel_biIdeal, EnvelopeProperty.ProjSubvariety.BlowupCenter.multiModelAnHomeomorph, EnvelopeProperty.ProjSubvariety.BlowupCenter.isPSH_multiModel_comp_iff, EnvelopeProperty.MultiModel.DegOnePres.sections_eq_piece_of_isProjNormal, EnvelopeProperty.MultiModel.DegOnePres.isPSH_comp_anHomeomorph_iff, EnvelopeProperty.MultiModel.DegOnePres.hasEnvelopeProperty_toProjSubvariety_iff, EnvelopeProperty.MultiModel.isIntegral_monomial_of_isAdmissible_of_le, EnvelopeProperty.MultiModel.sections_eq_piece_of_le_coordAlg}
\leanok
\uses{lem:gl-bridge-points, lem:gl-sections, def:gl-psh, def:pg-subvariety-psh, def:pg-blowup-psh, thm:pg-psh2, lem:gl-model-change}
\begin{enumerate}
\item In Lemma~\ref{lem:gl-bridge-points}(1), if $S$ is integrally closed (i.e.\ $V(I)$ is projectively normal), then $R_a=S_a$ for $a\ge0$,
$\PSH(\mathcal X,t)\circ\tau=\PSH(V(I),t)$ for $t\ge0$, and $\mathrm{EP}_{\mathcal X}(t)\iff\mathrm{EP}_{V(I)}(t)$ for $t\ge0$. (For $a<0$, $R_a$ need not
vanish: if $V(I)$ is a point, $S=k[w_0T]$ and $R_a=k\,w_0^a$ for every $a\in\Z$, so constants are $t$-psh on $\mathcal X$ for $t<0$ while
$\PSH(V(I),t)=\varnothing$.)
\item In Lemma~\ref{lem:gl-bridge-points}(2), $\PSH(\mathcal X,\theta)\circ\vartheta^{-1}=\PSH(B,\theta)$ for $\theta\in\R^2_{\ge0}$; for every $\theta$,
$\PSH(B,\theta)\circ\vartheta\subseteq\PSH(\mathcal X,\theta)$ and $\PSH(\mathcal X,\theta)\circ\vartheta^{-1}\subseteq\PSH(B,\max(\theta,0))$. The model of a blowup
$B$ is the quotient model of its ideal (Remark~\ref{rmk:gl-model-examples}(1)). The bound $\theta\ge0$ cannot be dropped: $\PSH(B,\theta)=\varnothing$
unless $\theta\ge0$, while model sections exist in degrees outside $\N^2$ (for the blowup of $\PP^2$ at $[1:0:0]$, $x_1/x_0\in R_{(2,-1)}$).
\end{enumerate}
\end{lemma}

\begin{proof}
Let $f\in R_a$ with admissible $b$, and let $I_b\subseteq S$ be the ideal generated by $\tilde z^\beta T^b$, $\beta\in\mathrm{Mon}(b)$. If $S$ is
generated by the coordinates (in (1), first change the coordinates to $w$, which does not change $R$, Lemma~\ref{lem:gl-bridge-points}(1)), $S_{a+b}\subseteq I_b$, so $fT^aI_b\subseteq I_b$ and $fT^a$ is integral over $S$ (determinant trick; $I_b$ is finitely
generated and faithful). (1) If $S$ is integrally closed, $fT^a\in S_a$ for $a\ge0$. FS functions of classes $a<0$ are replaced by classes $\max(a,0)$; the FS families of nonnegative classes then coincide,
hence the psh classes and the envelope property. (2) Directly from $fT^aI_b\subseteq I_b$, $(fT^a)^nI_b\subseteq I_b\subseteq S$: the admissible degree $b$ works for every power (no
finiteness of $S[fT^a]$ is needed). Hence for a model FS function $\varphi=m^{-1}\log\max_j|f_j|$ of class $a$, and every
$n$, $\varphi=(mn)^{-1}\log\max_{j,\beta}|f_j^n\tilde z^\beta T^{na+b}|$ is a Fubini--Study function on $B$ of class $a+b/(mn)$; by PSH2 on
$B$ (Theorem~\ref{thm:pg-psh2}) $\varphi\in\PSH(B,a)$. A model FS function of a class $a\in\Q^2$ is one of class $a^+=\max(a,0)$ (multiply by the
monomials of degree $a^+-a$, whose maximum is $1$), so a decreasing limit of model FS functions with classes $a_i\to\theta\ge0$ is a decreasing limit of
functions in $\PSH(B,a_i^+)$ with $a_i^+\to\theta$, and lies in $\PSH(B,\theta)$ by PSH2. Conversely polynomial sections are model sections.
\end{proof}

\begin{definition}[Smooth models]\label{def:gl-smooth}
\lean{EnvelopeProperty.MultiModel.IsSmooth, EnvelopeProperty.MultiModel.isSmooth_iff_of_isNormal}
\leanok
\uses{lem:gl-charts}
A model is \emph{smooth} if every chart algebra $A_j$ of Lemma~\ref{lem:gl-charts} is a smooth $k$-algebra and integrally closed.
\end{definition}

\begin{theorem}[Smooth case; interface to the smooth chapter]\label{thm:gl-smooth-ep}
\lean{EnvelopeProperty.MultiModel.hasEnvelopeProperty_of_isSmooth}
\leanok
\uses{def:gl-smooth, def:gl-psh, thm:pn-ep, cor:sg-smooth-model}
\source{fhx-envelope-20261028:thm:smooth-EP}
Let $\mathcal Y$ be a smooth model of rank $1$ and $t\ge0$. Then $\mathrm{EP}_{\mathcal Y}(t)$.
\end{theorem}

\begin{proof}
This is \cite[Theorem~3.1]{FHX} for the very ample class $\mathcal O(1)$ and its multiples, produced by the chapter on general smooth varieties
(Corollary~\ref{cor:sg-smooth-model}; for $\mathcal Y=\PP^n$ it is Theorem~\ref{thm:pn-ep}). By Lemma~\ref{lem:gl-ep-basic}(1),(2) it suffices that $\mathrm{EP}(1)$ holds.
\end{proof}

\begin{corollary}[Smooth bigraded models]\label{cor:gl-smooth-ep-nef}
\lean{EnvelopeProperty.MultiModel.hasEnvelopeProperty_of_isSmooth_of_nonneg, EnvelopeProperty.MultiModel.IsSmooth.veronese, EnvelopeProperty.MultiModel.IsSmooth.normalization}
\leanok
\uses{thm:gl-smooth-ep, lem:gl-veronese, lem:gl-ep-basic, def:gl-smooth}
Let $\mathcal Y$ be a smooth model of rank $r$ and $\theta\in\R^r_{\ge0}$. Then $\mathrm{EP}_{\mathcal Y}(\theta)$.
\end{corollary}

\begin{proof}
By Lemma~\ref{lem:gl-ep-basic}(1) with rational positive $\theta_j\downarrow\theta$, and (2), it suffices to treat positive $a\in\N^r$. The chart algebras of
$\mathcal Y_{(a)}$ at the coordinates $\tilde z^\beta U$ are the localizations of the chart algebras of $\mathcal Y$ at a chart $j$ in the support of $\beta$ (every
block occurs in $\beta$ as $a$ is positive) at the element $\tilde z^\beta/\tilde z_j^a$, hence smooth and integrally closed. Apply
Lemma~\ref{lem:gl-veronese}, Remark~\ref{rmk:gl-bridges} and Theorem~\ref{thm:gl-smooth-ep}.
\end{proof}

\subsection{Homogeneous functions}\label{subsec:gl-homogeneous}

\begin{definition}[Scaling]\label{def:gl-scaling}
\lean{EnvelopeProperty.MultiModel.rpow, EnvelopeProperty.MultiModel.rpowHomeomorph, EnvelopeProperty.MultiModel.scaleFun, EnvelopeProperty.MultiModel.IsScaleHomogeneous, EnvelopeProperty.MultiModel.homogenization, EnvelopeProperty.MultiModel.legendreSlice}
\leanok
\uses{def:gl-points}
\source{fhx-envelope-20261028:eq:scaling-action, bj22-trivval:sec:homogpsh}
For $t>0$ and $v\in\mathcal X^\an$, $|p|_{t\cdot v}=|p|_v^t$ defines $t\cdot v\in\mathcal X^\an$; on functions $(t\cdot u)(v)=t\,u(t^{-1}\cdot v)$; $u$ is
\emph{homogeneous} if $t\cdot u=u$ for all $t>0$. The \emph{homogenization} is $\hat u=\inf_{t>0}t\cdot u$, and the \emph{Legendre slice} at
$\lambda\in\R$ is $\hat u_\lambda=\widehat{u-\lambda}$.
\end{definition}

\begin{lemma}[Scaling]\label{lem:gl-scaling}
\lean{EnvelopeProperty.MultiModel.continuous_rpow_prod, EnvelopeProperty.MultiModel.rpow_triv, EnvelopeProperty.MultiModel.logSecAbs_rpow, EnvelopeProperty.MultiModel.FSData.scaleFun_eval, EnvelopeProperty.MultiModel.isPSH_scaleFun_iff, EnvelopeProperty.MultiModel.isHomFS_iff, EnvelopeProperty.MultiModel.IsPSH.apply_triv_eq_zero, EnvelopeProperty.MultiModel.IsPSH.nonpos, EnvelopeProperty.MultiModel.homogenization_le, EnvelopeProperty.MultiModel.isScaleHomogeneous_homogenization, EnvelopeProperty.MultiModel.IsScaleHomogeneous.le_homogenization, EnvelopeProperty.MultiModel.homogenization_iInf, EnvelopeProperty.MultiModel.scaleFun_iSup, EnvelopeProperty.MultiModel.scaleFun_uscReg, EnvelopeProperty.MultiModel.IsScaleHomogeneous.uscReg_iSup}
\leanok
\uses{def:gl-scaling, def:gl-psh, thm:gl-fs-system-theory}
\source{bj22-trivval:sec:homogpsh, bj22-trivval:lem:homogdec, bj22-trivval:lem:genpos, fhx-envelope-20261028:prop:BJ-package}
(1) $(t,v)\mapsto t\cdot v$ is continuous, $\star$ is fixed, and $\log|f|(t\cdot v)=t\log|f|(v)$. (2) For $\varphi=m^{-1}\max_j\{\log|f_j|+\lambda_j\}$,
$t\cdot\varphi=m^{-1}\max_j\{\log|f_j|+t\lambda_j\}$; hence $t\cdot\PSH(\theta)=\PSH(\theta)$ and homogeneous FS functions are those of
$\mathcal H_{\mathrm{hom}}$. (3) Homogeneous psh functions vanish at $\star$ and are $\le0$. (4) If $u_i\downarrow u$ then $\hat u_i\downarrow\hat u$;
$\hat u\le u$ is the largest homogeneous function below $u$. (5) Scaling commutes with pointwise suprema and usc regularization; in particular
$(\sup_iu_i)^\star$ is homogeneous if all $u_i$ are.
\end{lemma}

\begin{proof}
(1)--(2) are the formulas. A decreasing net of FS functions for $u$ gives one for $t\cdot u$ with the same classes. (3) $u(\star)=(t\cdot u)(\star)=tu(\star)$
for all $t$, and $u(\star)\in\R$; then $u\le u(\star)=0$. (4) is \cite[Lemma~6.13]{BJ22}: $\hat u\le\hat u_i\le u_i$, so $\lim\hat u_i$ is homogeneous between
$\hat u$ and $u$. (5) $v\mapsto t\cdot v$ is a homeomorphism.
\end{proof}

\begin{lemma}[A one-variable Legendre formula]\label{lem:gl-legendre}
\lean{EnvelopeProperty.legendreMax, EnvelopeProperty.iInf_sup_add_mul_eq_legendreMax, EnvelopeProperty.legendre_weights_pos}
\leanok
\source{bj22-trivval:lem:legendre}
Let $\lambda_1,\dots,\lambda_n\in\Q$ with $\max_i\lambda_i\ge0$ and $y\in[-\infty,\infty)^n$. Then
\[
\inf_{t>0}\max_i\{y_i+t\lambda_i\}=\max\Bigl\{\max_{\lambda_i\ge0}y_i,\ \max_{\lambda_i>0>\lambda_l}\frac{\lambda_iy_l-\lambda_ly_i}{\lambda_i-\lambda_l}\Bigr\},
\]
with $0\cdot(-\infty)=0$; the coefficients $\lambda_i/(\lambda_i-\lambda_l)$, $-\lambda_l/(\lambda_i-\lambda_l)$ are positive rationals with sum $1$.
\end{lemma}

\begin{proof}
Let $h(t)=\max_i\{y_i+t\lambda_i\}$ (indices with $y_i=-\infty$ omitted; if all are omitted both sides are $-\infty$). The right side $\rho$ is $\le h(t)$ for
every $t\ge0$: $y_i+t\lambda_i\ge y_i$ for $\lambda_i\ge0$, and for $\lambda_i>0>\lambda_l$ the convex combination with the displayed weights of $y_i+t\lambda_i$ and
$y_l+t\lambda_l$ equals $(\lambda_iy_l-\lambda_ly_i)/(\lambda_i-\lambda_l)$. As $h$ is continuous on $[0,\infty)$, $\inf_{t>0}h=\inf_{t\ge0}h$. If all remaining $\lambda_i<0$,
$h\to-\infty$ and $\rho=-\infty$ (every term involves an omitted index with positive weight). Otherwise $h$ is convex piecewise linear and eventually
nondecreasing, so it attains its minimum at some $t^*\ge0$. If $t^*=0$, the right derivative $\max\{\lambda_i:i\text{ active at }0\}\ge0$, so an active $i$ has
$\lambda_i\ge0$ and $h(0)=y_i\le\rho$. If $t^*>0$, the active set contains $i$ with $\lambda_i\ge0$ and $l$ with $\lambda_l\le0$; if one of them is $0$, $h(t^*)=y_\cdot\le\rho$;
otherwise $\lambda_i>0>\lambda_l$ and $y_i+t^*\lambda_i=y_l+t^*\lambda_l=h(t^*)$, so $h(t^*)$ is the displayed convex combination.
\end{proof}

\begin{theorem}[Homogenization of Fubini--Study functions]\label{thm:gl-homog-fs}
\lean{EnvelopeProperty.MultiModel.IsRatFS, EnvelopeProperty.MultiModel.IsRatFS.isHomFS_homogenization, EnvelopeProperty.MultiModel.IsFS.homogenization_eq_bot}
\leanok
\uses{lem:gl-legendre, lem:gl-scaling, def:gl-fs}
\source{bj22-trivval:thm:homogFS}
Let $\varphi=m^{-1}\max_t\{\log|f_t|+\lambda_t\}\in\mathcal H(\mathcal X,a)$ with $\lambda_t\in\Q$. If $\varphi(\star)\ge0$, then
$\hat\varphi\in\mathcal H_{\mathrm{hom}}(\mathcal X,a)$; otherwise $\hat\varphi\equiv-\infty$.
\end{theorem}

\begin{proof}
$\hat\varphi(v)=m^{-1}\inf_{s>0}\max_t\{\log|f_t|(v)+s\lambda_t\}$ by Lemma~\ref{lem:gl-scaling}(2). If $\varphi(\star)=m^{-1}\max\lambda_t<0$ the infimum is
$-\infty$. Otherwise Lemma~\ref{lem:gl-legendre} (with $y_t=\log|f_t|(v)$, indices with $f_t=0$ omitted) writes $m\hat\varphi$ as a maximum of
$\log|f_t|$ ($\lambda_t\ge0$) and of $w\log|f_t|+w'\log|f_l|$ with positive rational $w+w'=1$. With $b\in\mathbf Z_{>0}$ clearing denominators these are
$b^{-1}\log|f_t^b|$ and $b^{-1}\log|f_t^{bw}f_l^{bw'}|$, with sections in $R_{bma}$ (Lemma~\ref{lem:gl-sections}(2)), so
$\hat\varphi\in\mathcal H_{\mathrm{hom}}(\mathcal X,a)$, nonzero since some $\lambda_t\ge0$ gives a nonzero section.
\end{proof}

\begin{corollary}[Homogeneous psh functions]\label{cor:gl-pshhom}
\lean{EnvelopeProperty.MultiModel.IsPSH.exists_isRatFS_net, EnvelopeProperty.MultiModel.IsPSH.exists_isHomFS_net, EnvelopeProperty.MultiModel.IsPSH.isPSHHom_homogenization, EnvelopeProperty.MultiModel.IsPSH.homogenization_eq_bot}
\leanok
\uses{thm:gl-homog-fs, lem:gl-scaling, def:gl-psh}
\source{bj22-trivval:cor:pshhom, bj22-trivval:lem:genpos}
Let $u\in\PSH(\mathcal X,\theta)$ with $u(\star)\ge0$. Then $\hat u\in\PSH_{\mathrm{hom}}(\mathcal X,\theta)$, and $\hat u$ is the limit of a decreasing net in
$\mathcal H_{\mathrm{hom}}(\mathcal X,a_\alpha)$ with $a_\alpha\to\theta$. If $u(\star)<0$, $\hat u\equiv-\infty$.
\end{corollary}

\begin{proof}
Let $\varphi_\alpha\in\mathcal H(a_\alpha)$ decrease to $u$. Replace real weights by rational ones: with dyadic roundings $\lambda^{(n)}\ge\lambda$ within $2^{-n}$,
the functions $\varphi_{\alpha,n}+2^{1-n}$ form a decreasing net with limit $u$ for the order $(\alpha',n')\ge(\alpha,n)$ iff $\alpha'\ge\alpha$ and $n'\ge n+2$ or
$(\alpha',n')=(\alpha,n)$ (since $\varphi_{\alpha,n}\in[\varphi_\alpha,\varphi_\alpha+2^{-n}]$); this uses no diagonal lemma. Then $\varphi_\alpha(\star)\ge u(\star)\ge0$, the
$\hat\varphi_\alpha\in\mathcal H_{\mathrm{hom}}(a_\alpha)$ (Theorem~\ref{thm:gl-homog-fs}) decrease to $\hat u$ (Lemma~\ref{lem:gl-scaling}(4)), and
$\hat u(\star)=0$, so $\hat u\in\PSH_{\mathrm{hom}}(\theta)$ by definition. If $u(\star)<0$, $\hat u\le\widehat{u(\star)}=-\infty$.
\end{proof}

\begin{lemma}[Convexity along rays]\label{lem:gl-action-convex}
\lean{EnvelopeProperty.MultiModel.rayFun, EnvelopeProperty.MultiModel.IsPSH.scaleFun_le, EnvelopeProperty.MultiModel.IsPSH.scaleFun_convex, EnvelopeProperty.MultiModel.IsPSH.scaleFun_eq_bot_or_ne_bot, EnvelopeProperty.MultiModel.IsPSH.convexOn_rayFun}
\leanok
\uses{lem:gl-scaling, def:gl-psh}
\source{bj22-trivval:lem:actionconvex, bj22-trivval:prop:pshconvex}
For $u\in\PSH(\mathcal X,\theta)$ and $v\in\mathcal X^\an$, the function $g(t)=(t\cdot u)(v)$ on $(0,\infty)$ is either identically $-\infty$ or finite and
convex, and $g(t)\le t\,u(\star)$.
\end{lemma}

\begin{proof}
For an FS function, $g(t)=m^{-1}\max_j\{\log|f_j|(v)+t\lambda_j\}$ by Lemma~\ref{lem:gl-scaling}(2), with the terms with $|f_j|_v=0$ equal to
$-\infty$ for all $t$; so $g$ is $\equiv-\infty$ or a finite maximum of affine functions. For $u=\lim\varphi_\alpha$ the functions $g_\alpha$ decrease to $g$, so $g$ is
convex with values in $[-\infty,\infty)$ (the convexity inequality passes to the limit of a decreasing net). If $g(t_1)=-\infty$ and $t_0\ne t_1$, choose $t_2$ with $t_0$
strictly between $t_1$ and $t_2$; then $g(t_0)\le\lambda g(t_1)+(1-\lambda)g(t_2)=-\infty$. So $g\equiv-\infty$ or $g$ is finite. The bound is
$u(t^{-1}\cdot v)\le u(\star)$.
\end{proof}

\begin{theorem}[Homogeneous decomposition]\label{thm:gl-homog-decomposition}
\lean{EnvelopeProperty.MultiModel.IsPSH.isPSHHom_legendreSlice, EnvelopeProperty.MultiModel.legendreSlice_antitone, EnvelopeProperty.MultiModel.legendreSlice_mono, EnvelopeProperty.MultiModel.IsPSH.eq_iSup_legendreSlice}
\leanok
\uses{lem:gl-action-convex, cor:gl-pshhom}
\source{bj22-trivval:thm:homogpsh, bj22-trivval:sec:homogdecomp, fhx-envelope-20261028:eq:Legendre-reconstruction}
Let $u\in\PSH(\mathcal X,\theta)$ and $b=u(\star)$. For $\lambda<b$, $\hat u_\lambda\in\PSH_{\mathrm{hom}}(\mathcal X,\theta)$; $\lambda\mapsto\hat u_\lambda$ is
decreasing; $u\mapsto\hat u_\lambda$ is increasing; and $u=\sup_{\lambda<b}(\hat u_\lambda+\lambda)$ pointwise on $\mathcal X^\an$.
\end{theorem}

\begin{proof}
$\hat u_\lambda=\widehat{u-\lambda}$ and $(u-\lambda)(\star)>0$, so Corollary~\ref{cor:gl-pshhom} applies; monotonicity in $\lambda$ and $u$ is clear from the definition.
Fix $v$ and $g(t)=(t\cdot u)(v)$. If $g\equiv-\infty$, then $u(v)=-\infty$ and all $\hat u_\lambda(v)=-\infty$. Otherwise $g$ is finite convex on $(0,\infty)$
(Lemma~\ref{lem:gl-action-convex}), hence continuous, and Fenchel duality on the open interval gives $g(1)=\sup_\lambda\{\lambda-g^*(\lambda)\}$ with
$g^*(\lambda)=\sup_{t>0}\{t\lambda-g(t)\}=-\hat u_\lambda(v)$. For $\lambda>b$, $g^*(\lambda)\ge\sup_tt(\lambda-b)=\infty$; the concave function $\lambda\mapsto\hat u_\lambda(v)+\lambda$ satisfies
$\hat u_b(v)+b\le\lim_{\lambda\uparrow b}(\hat u_\lambda(v)+\lambda)$, so the supremum may be taken over $\lambda<b$.
\end{proof}

\subsection{Divisorial points and the comparison theorem}\label{subsec:gl-comparison}

\begin{definition}[Gauss extension and divisorial points]\label{def:gl-div}
\lean{EnvelopeProperty.MultiModel.gaussPt, EnvelopeProperty.MultiModel.IsDivisorial, EnvelopeProperty.MultiModel.isDivisorial_of_val_eq}
\leanok
\uses{def:gl-points, def:pn-gauss, def:pn-div}
\source{fhx-envelope-20261028:subsec:prelim-berkovich, bj22-trivval:sec:Gauss}
For $v\in\mathcal X^\an$, $\sigma(v)$ is the seminorm $|\sum_\lambda p_\lambda t^\lambda|_{\sigma(v)}=\max_\lambda|p_\lambda|_ve^{-\lambda}$ on $S[t]$. $v$ is
\emph{divisorial} if there are $c\in\Q_{>0}$ and a discrete valuation ring $S[t]\subseteq R'\subseteq\operatorname{Frac}(S[t])$, essentially of finite type
over $k$, with $\log|g|_{\sigma(v)}=-c\operatorname{ord}_{R'}(g)$ for $g\in S[t]$. $\mathcal X^{\divv}$ is the set of divisorial points.
\end{definition}

\begin{lemma}[Divisorial points]\label{lem:gl-div-generic}
\lean{EnvelopeProperty.MultiModel.isDivisorial_triv, EnvelopeProperty.MultiModel.IsDivisorial.hasGenericSupport, EnvelopeProperty.MultiModel.IsDivisorial.mem_valSet, EnvelopeProperty.MultiModel.IsDivisorial.secAbs_ne_zero, EnvelopeProperty.MultiModel.IsDivisorial.logSecAbs_ne_bot}
\leanok
\uses{def:gl-div, lem:gl-sections}
$\star\in\mathcal X^{\divv}\subseteq\mathcal X^{\mathrm{val}}$; in particular every nonzero section has finite logarithm at every divisorial point.
\end{lemma}

\begin{proof}
For $\star$ take $R'=S[t]_{(t)}$, a discrete valuation ring because $(t)$ is a principal prime of height one ($S[t]/(t)=S$ is a domain), with
$\operatorname{ord}(\sum p_\lambda t^\lambda)=\min\{\lambda:p_\lambda\ne0\}$, and $c=1$. A divisorial point is nonzero on $S\setminus0$ because $\operatorname{ord}_{R'}$ is finite
there; sections have finite logarithm at points of $\mathcal X^{\mathrm{val}}$ by Lemma~\ref{lem:gl-sections}(3).
\end{proof}

\begin{lemma}[Points attached to Rees valuations]\label{lem:gl-rees-points}
\lean{EnvelopeProperty.MultiModel.chartDivPoint, EnvelopeProperty.MultiModel.chartDivPoint_val_ofPiece, EnvelopeProperty.MultiModel.isDivisorial_chartDivPoint, EnvelopeProperty.MultiModel.gaussPt_chartDivPoint_val, EnvelopeProperty.MultiModel.chartGaussDVR_ord_of_eq_monomial, EnvelopeProperty.MultiModel.reesDVR, EnvelopeProperty.MultiModel.reesDVR_torus, EnvelopeProperty.MultiModel.isDivisorial_reesPoint, EnvelopeProperty.BirationalDVR.ofMvGauss}
\leanok
\uses{lem:gl-charts, def:gl-div, def:bj-rees, lem:bj-rees-invariant, lem:bj-gauss-dvr, lem:gl-multigraded-extension}
\source{bj22-trivval:defi:Rees, bj22-trivval:lem:Shilov}
Let $j$ be a chart, $B=A_j[t]$, and $I\subseteq B$ generated by $g_\ell=h_\ell t^{e_\ell}$ ($0\ne h_\ell\in A_j$) with $t^N\in I$. For every
$V\in\mathrm{RV}(I;g)$, $b=\operatorname{ord}_V(t)>0$, and $|pT^a|_{v_V}=\exp(-\operatorname{ord}_V(p/\tilde z_j^a)/b)$ ($p\in S_a$) defines a divisorial point
$v_V\in U_j$ with $\log|g|_{\sigma(v_V)}=-\operatorname{ord}_V(g^{(j)})/b$, where $g^{(j)}\in B$ is the dehomogenization of $g\in S[t]$ homogeneous in $T$.
\end{lemma}

\begin{proof}
The proof of Lemma~\ref{lem:bj-rees-invariant} uses only that $B=A[t]$ for a domain $A$ and that the generators are monomials in $t$ with $t^N\in I$;
it gives $b>0$ and $\operatorname{ord}_V(\sum f_\lambda t^\lambda)=\min_\lambda(\operatorname{ord}_Vf_\lambda+\lambda b)$. The coordinate ratios lie in $A_j\subseteq V$ and equal $1$ at
$j_c$, so Lemma~\ref{lem:gl-multigraded-extension} gives $v_V$, with $|z^{(c)}_{j_c}|=1$. Let $R'$ be the Gauss extension of $V$ in the $r$ variables
$z^{(c)}_{j_c}$ (Lemma~\ref{lem:bj-gauss-dvr}, iterated), a discrete valuation ring essentially of finite type, with $K(T)(t)=\operatorname{Frac}(V)(z_j)$; it
contains $S[t]$, since $pT^at^\lambda=(p/\tilde z_j^a)t^\lambda\cdot z_j^a$. For $g=\sum_{a,\lambda}p_{a,\lambda}T^at^\lambda$,
$\operatorname{ord}_{R'}(g)=\min_{a,\lambda}(\operatorname{ord}_V(p_{a,\lambda}/\tilde z_j^a)+\lambda b)=-b\log|g|_{\sigma(v_V)}$; take $c=1/b$.
\end{proof}

\begin{lemma}[Shilov boundary]\label{lem:gl-shilov}
\lean{EnvelopeProperty.MultiModel.IsFinFS.exists_finset_le_of_le, EnvelopeProperty.MultiModel.IsFinFS.exists_finset_iSup_sub_eq, EnvelopeProperty.MultiModel.IsFinFS.exists_shilovSup}
\leanok
\uses{lem:gl-rees-points, thm:bj-rees, lem:gl-cones, lem:gl-sections, lem:bj-shilov, lem:bj-finite-fs}
\source{bj22-trivval:lem:Shilov, bj22-trivval:lem:supPL}
For every $\psi\in\mathcal H_{\mathrm{fin}}(\mathcal X,b)$ there is a finite set $\Sigma\subseteq\mathcal X^{\divv}$ containing $\star$ such that
$\sup_{\mathcal X^\an}(\varphi-\psi)=\max_\Sigma(\varphi-\psi)$ for every $a\in\Q^r$ and every $\varphi\in\mathcal H(\mathcal X,a)$ with rational weights.
\end{lemma}

\begin{proof}
As Lemma~\ref{lem:bj-shilov}. After subtracting a rational constant and passing to powers, $\psi=m^{-1}\max_i(\log|f_i|+\lambda_i)$ with
$f_i\in R_{mb}$ without common zero, $\lambda_i\in\mathbf Z_{\le0}$, $\max\lambda_i=0$. For each chart $j$, $g_i=f_i/\tilde z_j^{mb}\in A_j$ generate the unit
ideal (proof of Lemma~\ref{lem:gl-cones}(2)), so $I^{(j)}=(g_it^{-\lambda_i})\subseteq A_j[t]$ contains $t^N$, $N=\max(-\lambda_i)$. Let $\Sigma$ consist of $\star$
and the points $v_V$, $V\in\mathrm{RV}(I^{(j)})$, over all charts (Lemma~\ref{lem:gl-rees-points}). Write $\varphi=M^{-1}\max_\ell(\log|\sigma_\ell|+\mu_\ell)$ with
$\sigma_\ell\in R_{Ma}$, $\mu_\ell\in\mathbf Z$, put $c=\max_\Sigma(\varphi-\psi)\in\Q$ (finite since $\varphi(\star)$ is finite and $\psi(\star)=0$), choose $r$ with
$rMc\in\mathbf Z$, and $F_\ell=\sigma_\ell^{rm}t^{a_\ell}$ with $a_\ell=m(rMc-r\mu_\ell)\in\N$. At any $v$, $(\varphi-\psi)(v)\le c$ iff
$\log|F_\ell|_{\sigma(v)}\le rM\max_i(\log|f_i|_v+\lambda_i)$ for all $\ell$. On a chart $j$ with $|z^{(c)}_{j_c}|_v=1$ this reads
$|F^{(j)}_\ell|_w\le\max_i|g_it^{-\lambda_i}|_w^{rM}$ for the multiplicative seminorm $w=\sigma(\iota_j(v))$ on $A_j[t]$, which is $\le1$ on $A_j[t]$ and has
$|t^N|_w>0$; at $v=v_V$ it reads $\operatorname{ord}_V(F^{(j)}_\ell)\ge rM\operatorname{ord}_V(I^{(j)})$. Theorem~\ref{thm:bj-rees} passes from the Rees points to all $v$.
\end{proof}

\begin{lemma}[PL functions]\label{lem:gl-pl}
\lean{EnvelopeProperty.MultiModel.plSystem, EnvelopeProperty.MultiModel.IsPL, EnvelopeProperty.MultiModel.isPL_iff, EnvelopeProperty.MultiModel.separatesPoints_isPL, EnvelopeProperty.MultiModel.dense_setOf_isPL, EnvelopeProperty.MultiModel.exists_antitone_isPL_net, EnvelopeProperty.MultiModel.le_of_forall_isPL_gt, EnvelopeProperty.FSSystemC.toFSSystemNonneg}
\leanok
\uses{lem:gl-fs-system, lem:bj-pl-lattice, thm:bj-pl-dense, lem:bj-usc-pl}
\source{bj22-trivval:lem:FSdense, bj22-trivval:thm:PLdense}
Let $\mathrm{PL}(\mathcal X)$ be the set of differences of functions in $\bigcup_{a\in\Q^r_{\ge0}}\mathcal H_{\mathrm{fin}}(\mathcal X,a)$. It is a $\Q$-vector space
of continuous real functions, contains $\Q$, is stable under finite maxima and minima, separates points, and is dense in $C^0(\mathcal X^\an)$. For every usc
$h$, the set $\{\psi\in\mathrm{PL}:\psi>h\}$ is nonempty, directed downwards, with infimum $h$.
\end{lemma}

\begin{proof}
As Lemmas~\ref{lem:bj-pl-lattice}, \ref{lem:bj-usc-pl} and Theorem~\ref{thm:bj-pl-dense}: common classes by (FS2) with coordinatewise maxima; points
$v\ne w$ differ on a homogeneous $p\in S_a$ (both are multigraded), and $\max\{\log|p|,-N\}\in\mathcal H_{\mathrm{fin}}(a)$ separates them for $N$ large; the
lattice Stone--Weierstrass theorem on the compact space $\mathcal X^\an$; Urysohn for the last claim.
\end{proof}

\begin{theorem}[Divisorial comparison]\label{thm:gl-comparison}
\lean{EnvelopeProperty.MultiModel.IsPSH.le_of_forall_isDivisorial_le, EnvelopeProperty.MultiModel.IsPSH.eq_of_forall_isDivisorial_eq, EnvelopeProperty.MultiModel.IsPL.exists_finset_iSup_sub_eq, EnvelopeProperty.MultiModel.IsPSH.le_of_forall_isDivisorial_le_isPL}
\leanok
\uses{lem:gl-shilov, lem:gl-pl, cor:gl-pshhom, thm:bj-divisorial-comparison, lem:gl-div-generic}
\source{bj22-trivval:thm:suppsh, bj22-trivval:cor:Hausdorff, fhx-envelope-20261028:prop:BJ-package}
Let $\theta\in\R^r$, $u\in\PSH(\mathcal X,\theta)$ and $h$ usc. If $u\le h$ on $\mathcal X^{\divv}$, then $u\le h$. Two $\theta$-psh functions (possibly of
different classes) that agree on $\mathcal X^{\divv}$ are equal.
\end{theorem}

\begin{proof}
Let $u_\alpha\in\mathcal H(a_\alpha)$ with rational weights decrease to $u$ (proof of Corollary~\ref{cor:gl-pshhom}). Let $\psi=\psi_1-\psi_2\in\mathrm{PL}$ with
$\psi>h$, and $\Sigma$ the set of Lemma~\ref{lem:gl-shilov} for $\psi_1$. Since $u_\alpha+\psi_2\in\mathcal H(a_\alpha+b_2)$,
$(u-\psi)(v)\le(u_\alpha+\psi_2-\psi_1)(v)\le\max_\Sigma(u_\alpha-\psi)$, and as $\Sigma$ is finite the right side tends to $\max_\Sigma(u-\psi)\le0$. So $u\le\psi$ for all
such $\psi$, and $u\le h$ by Lemma~\ref{lem:gl-pl}. For the last claim apply this in both directions.
\end{proof}

\begin{lemma}[Locality of usc regularization]\label{lem:gl-usc-local}
\lean{EnvelopeProperty.uscReg_restrict_apply, EnvelopeProperty.uscReg_add_continuous, EnvelopeProperty.MultiModel.IsPSH.le_of_forall_mem_le, EnvelopeProperty.MultiModel.IsPSH.eq_of_forall_mem_eq}
\leanok
\uses{thm:gl-comparison, def:gl-points}
Let $\Omega\subseteq\mathcal X^\an$ be open and $g$ bounded above. (1) $g^\star|_\Omega=(g|_\Omega)^\star$; if $g|_\Omega$ is usc, $g^\star=g$ on $\Omega$. (2) If
$e\colon\Omega\to\R$ is continuous, $(g+e)^\star=g^\star+e$ on $\Omega$. (3) If $\Omega\supseteq\mathcal X^{\divv}$, $u\in\PSH(\mathcal X,\theta)$ and $h$ is usc
with $u\le h$ on $\Omega$, then $u\le h$; psh functions agreeing on $\Omega$ are equal.
\end{lemma}

\begin{proof}
(1), (2): $g^\star(x)=\limsup_{y\to x}g(y)$ and neighbourhoods of $x\in\Omega$ may be taken inside $\Omega$. (3) is Theorem~\ref{thm:gl-comparison}.
\end{proof}

\subsection{Ample classes}\label{subsec:gl-ample}

\begin{proposition}[Envelopes of PL obstacles]\label{prop:gl-pl-envelope}
\lean{EnvelopeProperty.MultiModel.hasEnvelopeProperty_iff_forall_isPL, EnvelopeProperty.MultiModel.hasEnvelopeProperty_of_forall_isPL, EnvelopeProperty.MultiModel.HasEnvelopeProperty.isPSH_envelope, EnvelopeProperty.MultiModel.HasEnvelopeProperty.isPSH_envelope_of_continuous, EnvelopeProperty.FSSystemC.hasEnvelopeProperty_iff_forall_isPL}
\leanok
\uses{lem:gl-pl, thm:gl-fs-system-theory, prop:bj-envelope-criterion, def:gl-psh}
\source{fhx-envelope-20261028:prop:ample-envelopes, bj22-addendum:lem:contenv}
For $\theta\in N(\mathcal X)$: $\mathrm{EP}(\theta)$ holds iff $\mathrm P_\theta(\phi)\in\PSH(\mathcal X,\theta)$ for every $\phi\in\mathrm{PL}(\mathcal X)$.
\end{proposition}

\begin{proof}
As Proposition~\ref{prop:bj-envelope-criterion}: the family $\{\tau\le\phi\}$ contains the constant $\min\phi$ ($\theta\in N(\mathcal X)$,
Theorem~\ref{thm:gl-fs-system-theory}(2)); conversely $(\sup u_i)^\star$ is the infimum of the downward directed net $\mathrm P_\theta(\psi)$, $\psi\in\mathrm{PL}$,
$\psi>(\sup u_i)^\star$ (Lemma~\ref{lem:gl-pl}), and PSH2 at $\theta$ applies.
\end{proof}

\begin{lemma}[Fubini--Study envelopes]\label{lem:gl-fs-envelope}
\lean{EnvelopeProperty.MultiModel.envelope_eq_iSup_isFinFS, EnvelopeProperty.MultiModel.envelope_coe_eq_iSup_isFinFS, EnvelopeProperty.FSSystemC.envelope_eq_iSup_isFinFS}
\leanok
\uses{thm:gl-fs-system-theory, lem:bj-fs-envelope}
\source{bj22-addendum:lem:FSenv, bj22-trivval:lem:envlsc}
Let $a\in\Q^r$ be positive and $\phi$ bounded and lsc. Then $\mathrm P_a(\phi)=\sup\{\psi\in\mathcal H_{\mathrm{fin}}(\mathcal X,a):\psi\le\phi\}$.
\end{lemma}

\begin{proof}
As Lemma~\ref{lem:bj-fs-envelope}, with the regularization of Theorem~\ref{thm:gl-fs-system-theory}(2) ($\tau_\alpha\in\mathcal H_{\mathrm{fin}}(a_\alpha)$,
$a_\alpha\ge a$, $a_\alpha\to a$) and the rescaling by rational $\lambda_\alpha\le\min_ca_c/a_{\alpha,c}$, $\lambda_\alpha\to1$, which puts $\lambda_\alpha\tau_\alpha$ in
$\mathcal H_{\mathrm{fin}}(a)$ by (FS4), (FS2).
\end{proof}

\begin{remark}[Test configurations]\label{rmk:gl-test-configurations}
\uses{prop:gl-pl-envelope, lem:gl-fs-envelope, prop:gl-ample-rational}
\cite[Proposition~2.5(2)]{FHX} (base ideals of test configurations, \cite[Lemma~5.22]{BJ22}) is used in \cite{FHX} only in \S3 and in Step~2 of
\cite[Theorem~6.1]{FHX}. Here Step~2 uses Lemma~\ref{lem:gl-fs-envelope} instead (Proposition~\ref{prop:gl-ample-rational}), so test configurations and
\cite[Theorem~2.31]{BJ22} are not needed in this chapter.
\end{remark}

\begin{definition}[Fibres of finite morphisms]\label{def:gl-finite-fibres}
\lean{EnvelopeProperty.MultiModel.fieldPoint, EnvelopeProperty.MultiModel.fibreDeg, EnvelopeProperty.MultiModel.FiniteHom.fibre, EnvelopeProperty.MultiModel.FiniteHom.fibreDeg_eq_finrank}
\leanok
\uses{def:gl-maps, lem:gl-pullback}
\source{fhx-envelope-20261028:sec:finite-descent}
Let $h\colon\mathcal T\to\mathcal X$ be a finite morphism of degree $d$ (Definition~\ref{def:gl-maps}) and $v\in\mathcal X^{\mathrm{val}}$, an absolute value on
$K$ (Definition~\ref{def:gl-points}) with completion $\mathcal H(v)$. Write $E\otimes_K\mathcal H(v)=\prod_wB_w$ (local Artinian factors) and
$d(w/v)=\dim_{\mathcal H(v)}B_w$.
\end{definition}

\begin{lemma}[Fibres over generic points]\label{lem:gl-finite-fibres}
\lean{EnvelopeProperty.MultiModel.FiniteHom.fibreEquiv, EnvelopeProperty.MultiModel.FiniteHom.fibreEquivFactors, EnvelopeProperty.MultiModel.FiniteHom.hasGenericSupport_of_anMap, EnvelopeProperty.MultiModel.FiniteHom.sum_fibreDeg, EnvelopeProperty.MultiModel.FiniteHom.fibre_triv, EnvelopeProperty.MultiModel.FiniteHom.exists_fieldPoint_eq, EnvelopeProperty.MultiModel.FiniteHom.eq_of_fieldPoint_eq, EnvelopeProperty.MultiModel.one_le_fibreDeg}
\leanok
\uses{def:gl-finite-fibres, lem:finite-flat-norm, lem:gl-charts, lem:gl-pullback}
\source{fhx-envelope-20261028:sec:finite-descent}
In Definition~\ref{def:gl-finite-fibres} the factors $B_w$ correspond bijectively to the points $w\in\mathcal T^\an$ with $h(w)=v$, these have generic support,
and $\sum_{w\mid v}d(w/v)=d$. The only point over $\star$ is $\star$.
\end{lemma}

\begin{proof}
On a chart $j$ with $|z_j|_v=1$, the chart algebra $B_j$ of $\mathcal T$ is finite over $A_j$ with $B_j\otimes_{A_j}K=E$ (localization), so
$B_j\otimes_{A_j}\mathcal H(v)=E\otimes_K\mathcal H(v)$; its maximal ideals are the points of $(\Spec B_j)^\an$ over $\iota_j(v)$ (\texttt{Berkovich/Finite},
Lemma~\ref{lem:finite-flat-norm}), and they have trivial kernel on $B_j$ because $E\otimes_K\mathcal H(v)$ is a $K$-algebra. Charts correspond since the
coordinates agree. $\sum d(w/v)=\dim_{\mathcal H(v)}E\otimes_K\mathcal H(v)=d$. The last claim is Lemma~\ref{lem:gl-pullback}.
\end{proof}

\begin{proposition}[Weighted coefficient norms]\label{prop:gl-weighted-norm}
\lean{EnvelopeProperty.MultiModel.FiniteHom.coeff_normForm_mem_sections, EnvelopeProperty.MultiModel.FiniteHom.avFSData, EnvelopeProperty.MultiModel.FiniteHom.isFS_avFSData, EnvelopeProperty.MultiModel.FiniteHom.eval_avFSData, EnvelopeProperty.MultiModel.FiniteHom.eval_avFSData_le, EnvelopeProperty.MultiModel.FiniteHom.eval_avFSData_eq, EnvelopeProperty.MultiModel.FiniteHom.eval_avFSData_triv, EnvelopeProperty.coeff_normForm_mem_of_isIntegral}
\leanok
\uses{lem:gl-finite-fibres, lem:gl-sections, thm:gl-comparison, lem:finite-flat-norm}
\source{fhx-envelope-20261028:prop:weighted-coefficient-norm, fhx-envelope-20261028:eq:divisorial-field-average}
Let $h\colon\mathcal T\to\mathcal X$ be finite of degree $d$ with $\mathcal X$ normal, $a\in\Q^r$, and
$\varphi=m^{-1}\max_j\{\log|s_j|+\lambda_j\}\in\mathcal H(\mathcal T,a)$, $s_j\in R_{ma}(\mathcal T)\subseteq E$. Write
$\Norm_{E(\mathbf T)/K(\mathbf T)}(\sum_j\mathbf T_js_j)=\sum_{|I|=d}c_I\mathbf T^I$. Then $c_I\in R_{dma}(\mathcal X)$, not all zero, and
\[
\mathrm{Av}_h(\varphi)\coloneqq(md)^{-1}\max_I\Bigl\{\log|c_I|+\sum_jI_j\lambda_j\Bigr\}\in\mathcal H(\mathcal X,a)
\]
satisfies $\mathrm{Av}_h(\varphi)(v)=d^{-1}\sum_{w\mid v}d(w/v)\varphi(w)$ for every $v\in\mathcal X^{\mathrm{val}}$. It does not depend on the presentation, and
$\varphi\le\varphi'$ on $\mathcal T^\an$ implies $\mathrm{Av}_h(\varphi)\le\mathrm{Av}_h(\varphi')$ on $\mathcal X^\an$, for $\varphi'$ of any class.
\end{proposition}

\begin{proof}
Let $b$ be admissible for all $s_j$ and $\beta\in\mathrm{Mon}(b)$. The linear form $\sum_j\mathbf T_js_j\tilde z^\beta T^{ma+b}$ has coefficients in $S_T$, hence is
integral over $S[\mathbf T]$; its norm from $E(T)(\mathbf T)$ to $K(T)(\mathbf T)$ is integral over $S[\mathbf T]$ and lies in $\operatorname{Frac}(S)(\mathbf T)$
($\operatorname{Frac}S=K(T)$ by (M3) and the coordinates), so it lies in $S[\mathbf T]$, $S$ being normal. It equals
$(\tilde z^\beta T^{ma+b})^d\sum_Ic_I\mathbf T^I$ (the norm of an element of the base field is its $d$-th power), so $c_I\tilde z^{d\beta}\in S_{d(ma+b)}$ for every
$\beta\in\mathrm{Mon}(b)$. The monomials $z^{d\beta}$ generate an ideal of $S$ containing $\mathrm{Mon}(b')$ for $b'\gg0$ (pigeonhole in each block), so $c_I\in
R_{dma}$. Some $c_I\ne0$ since the norm of a nonzero element is nonzero. Let $v\in\mathcal X^{\mathrm{val}}$, choose a chart with $|z_j|_v=1$; by
Lemma~\ref{lem:gl-sections}(4) everything is read on $A_j\subseteq B_j$ with $s_j/\tilde z_j^{ma}\in B_j$, and the identity
$\max_I\{\log|c_I|_v+\sum I_j\lambda_j\}=\sum_wd(w/v)\max_j\{\log|s_j|_w+\lambda_j\}$ is the norm identity for the weighted Gauss norm over
$E\otimes_K\mathcal H(v)$ (determinant on a local Artinian factor of length $\ell$ is the $\ell$-th power of the residue of the linear form, after extending
$\mathcal H(v)$ to a complete algebraically closed field; Lemma~\ref{lem:finite-flat-norm}, \texttt{Berkovich/NormIdentityGeneral}). Two presentations, or
$\varphi\le\varphi'$, give FS functions with equal (resp.\ ordered) values on $\mathcal X^{\mathrm{val}}\supseteq\mathcal X^{\divv}$; Theorem~\ref{thm:gl-comparison}
concludes, the functions being psh and continuous.
\end{proof}

\begin{theorem}[Finite descent]\label{thm:gl-finite-descent}
\lean{EnvelopeProperty.MultiModel.FiniteHom.hasEnvelopeProperty}
\leanok
\uses{prop:gl-weighted-norm, lem:gl-finite-fibres, thm:gl-comparison, def:gl-psh, lem:gl-pullback}
\source{fhx-envelope-20261028:thm:finite-normal-descent}
Let $h\colon\mathcal T\to\mathcal X$ be finite with $\mathcal X$ normal, and $\theta\in\R^r$. If $\mathrm{EP}_{\mathcal T}(\theta)$, then
$\mathrm{EP}_{\mathcal X}(\theta)$.
\end{theorem}

\begin{proof}
Let $(\varphi_i)$ be a nonempty family in $\PSH(\mathcal X,\theta)$ bounded above, $g=\sup_i\varphi_i$. The $\varphi_i\circ h\in\PSH(\mathcal T,\theta)$
(Lemma~\ref{lem:gl-pullback}), so $F\coloneqq(g\circ h)^\star\in\PSH(\mathcal T,\theta)$. Let $F_\alpha\in\mathcal H(\mathcal T,a_\alpha)$ decrease to $F$,
$a_\alpha\to\theta$. By Proposition~\ref{prop:gl-weighted-norm}, $G_\alpha=\mathrm{Av}_h(F_\alpha)\in\mathcal H(\mathcal X,a_\alpha)$ decrease on $\mathcal X^\an$, and
$G_\alpha(\star)=F_\alpha(\star)\to F(\star)\in\R$ (Lemma~\ref{lem:gl-finite-fibres}); so $G=\lim G_\alpha\in\PSH(\mathcal X,\theta)$ by definition. For $v\in\mathcal
X^{\mathrm{val}}$, finite sums of decreasing nets commute with limits, so $G(v)=d^{-1}\sum_{w\mid v}d(w/v)F(w)$. Since $g\circ h\le F\le g^\star\circ h$ (the
latter is usc and $\ge g\circ h$) and $h(w)=v$, $\varphi_i(v)\le g(v)\le G(v)\le g^\star(v)$ on $\mathcal X^{\mathrm{val}}$. By Theorem~\ref{thm:gl-comparison},
$\varphi_i\le G$ everywhere, so $g^\star\le G$, and $G\le g^\star$ everywhere. Thus $g^\star=G$ is $\theta$-psh.
\end{proof}

\begin{theorem}[Alterations]\label{thm:gl-dejong}
\lean{EnvelopeProperty.deJongAlterations, EnvelopeProperty.ProjSubvariety.alteration}
\leanok
\uses{rmk:gl-model-examples, def:gl-smooth}
\source{fhx-envelope-20261028:thm:normal-ample-EP, dejong96-alterations:page-0017}
Let $X=V(I)\subseteq\PP^N$ be a projective variety over $k$. There are $N'$ and a bihomogeneous prime $I_Y\subseteq k[x_0,\dots,x_N,y_0,\dots,y_{N'}]$, relevant
for both gradings, with $I_Y\cap k[x]=I$, such that the degree-zero fraction field $k(Y)$ of $S_Y=k[x,y]/I_Y$ is finite over $k(X)$ and the quotient
model $\mathcal Y$ of $I_Y$ (Remark~\ref{rmk:gl-model-examples}(1), rank $2$) is smooth (Definition~\ref{def:gl-smooth}).
\end{theorem}

\begin{proof}
Apply de Jong's Theorem~4.1 to the integral projective scheme
$X^s=\operatorname{Proj}(k[x]/I)$ and the proper closed subset $\emptyset$.
It gives a proper dominant generically finite map $f:Y\to X^s$ and an open
immersion $j:Y\to\overline Y$ into an integral regular projective scheme.
Since $X^s$ is proper over $k$, so is $Y$. As $\overline Y$ is separated over
$k$, the map $j$ is proper. Its image is therefore nonempty, open and closed
in the irreducible $\overline Y$, hence is all of $\overline Y$. Thus $j$ is
an isomorphism and $Y$ is regular and projective. The algebraically closed
field $k$ is perfect, so $Y$ is smooth over $k$.

Choose a projective embedding $\iota:Y\hookrightarrow\PP^{N'}$ and put
$L=k(Y)$. The two morphisms $Y\to\PP^N$ and $\iota$ give homogeneous generic
coordinates $a_i,b_m\in L$, with $a_{i_0}=b_{m_0}=1$ after normalization.
Let $I_Y$ be the kernel of the map
$k[x,y]\to L[T,S]$ sending $x_i$ to $a_iT$ and $y_m$ to $b_mS$.
This is a bihomogeneous prime ideal. The normalized coordinates show relevance
at $(i_0,m_0)$. Dominance of $f$ identifies the kernel in the $x$ block with
$I$, so $I_Y\cap k[x]=I$. The ratios of the $b_m$ generate $L$ as a field,
since $\iota$ is a closed immersion; the field generated by the $a_i$ is the
image of $k(X^s)$ in $L$. Consequently the degree-zero fraction field of the
quotient is $L$, finite over its $x$-ratio subfield by generic finiteness of $f$.
For any pair of nonvanishing generic coordinates, its chart algebra is
$k[a_i/a_j,b_m/b_n]$. This is the coordinate algebra of the affine open
$f^{-1}(D_+(x_j))\cap\iota^{-1}(D_+(y_n))$ in $Y$: the closed graph of $(f,\iota)$ meets the product of the two
standard affine charts in this affine scheme. These affine
algebras are smooth over $k$, and are integrally closed because $Y$ is regular
and integral. They give exactly the asserted smooth quotient model.
\end{proof}

\begin{remark}\label{rmk:gl-dejong}
\uses{thm:gl-dejong}
This is the sentence ``Choose an alteration $Y\to X$ with $Y$ smooth and projective \cite[Theorem~4.1]{deJong}'' of Step~1 of the proof of
\cite[Theorem~6.1]{FHX}, in the model: a projective $Y$ over $X$ embeds in $X\times\PP^{N'}$; $I_Y\cap k[x]=I$ says that $Y\to X$ is dominant, hence
surjective; finiteness of $k(Y)/k(X)$ says that it is generically finite. De Jong's theorem \cite[Theorem~4.1]{deJong} gives an alteration $\varphi_1\colon X_1\to X$
and an open immersion $j_1\colon X_1\to\bar X_1$ into a projective variety $\bar X_1$ that is a regular scheme (with a strict normal crossings divisor
containing $j_1(\varphi_1^{-1}(Z))$ and the boundary, not used here). For projective $X$ the map $\varphi_1$ is proper, so $j_1$ is also proper and
$X_1=\bar X_1$. Over the perfect field $k$, $\bar X_1$ is smooth \cite[Remark~4.2, 2.10]{deJong}, and regular rings are normal
(\cite[Tag~0567]{Stacks}). This is why Definition~\ref{def:gl-smooth} includes normality of the charts, which Mathlib does not prove for smooth algebras.
The model form is weaker than the source: it drops $Z$, the normal crossings divisor and generic \'etaleness. Checked against the source
(printed page 66). The scheme-theoretic alteration theorem and the conversion to these quotient models are both proved.
\end{remark}

\begin{definition}[Stein factorization in the model]\label{def:gl-stein}
\lean{EnvelopeProperty.ProjSubvariety.Alteration.baseModel, EnvelopeProperty.ProjSubvariety.Alteration.basePres, EnvelopeProperty.ProjSubvariety.Alteration.steinModel, EnvelopeProperty.ProjSubvariety.Alteration.steinHom}
\leanok
\uses{thm:gl-dejong, rmk:gl-model-examples, def:gl-maps}
\source{fhx-envelope-20261028:thm:normal-ample-EP}
In Theorem~\ref{thm:gl-dejong}, write $S_X\subseteq k(X)[T_1]$ and $S_Y\subseteq k(Y)[T_1,T_2]$ for the two quotient models (the $x$-coordinates of degree $e_1$,
the $y$-coordinates of degree $e_2$), so that $S_X\subseteq k(Y)[T_1]\subseteq k(Y)[T_1,T_2]$. Let $S_T$ be the integral closure of $S_X$ in $k(Y)[T_1]$,
$\mathcal T=(k(Y),S_T,x)$, and $h\colon\mathcal T\to\mathcal X$ the finite morphism (Remark~\ref{rmk:gl-model-examples}(3)).
\end{definition}

\begin{lemma}[Stein factorization]\label{lem:gl-stein}
\lean{EnvelopeProperty.ProjSubvariety.Alteration.isNormal_steinModel, EnvelopeProperty.ProjSubvariety.Alteration.isNormal_baseModel, EnvelopeProperty.ProjSubvariety.Alteration.hasEnvelopeProperty_iff_baseModel, EnvelopeProperty.ProjSubvariety.Alteration.steinBeta, EnvelopeProperty.ProjSubvariety.Alteration.steinHom_anMap_steinBeta, EnvelopeProperty.ProjSubvariety.Alteration.isPSH_comp_steinBeta, EnvelopeProperty.ProjSubvariety.Alteration.steinBeta_surjective, EnvelopeProperty.ProjSubvariety.Alteration.SteinBlowup, EnvelopeProperty.ProjSubvariety.Alteration.nonempty_steinBlowup, EnvelopeProperty.ProjSubvariety.Alteration.isPSH_steinTheta_iff, EnvelopeProperty.ProjSubvariety.Alteration.proj_steinTheta}
\leanok
\uses{def:gl-stein, lem:gl-model-change, lem:gl-veronese, rmk:gl-bridges, def:pg-blowup, lem:pg-blowup-basic, lem:pg-blowup-surjective, lem:wu-veronese, lem:wu-graded-closure, thm:bj-normalization-finite}
\source{fhx-envelope-20261028:thm:normal-ample-EP, fhx-envelope-20261028:cor:normal-weak-global}
In Definition~\ref{def:gl-stein}, assume $X$ projectively normal.
\begin{enumerate}
\item $\mathcal T$ is a normal model of rank $1$, finite over $\mathcal X$ of degree $[k(Y):k(X)]$, and $k(Y)$ is its function field.
\item $S_T\subseteq S_Y^\nu$, and $\beta(y)\coloneqq\tilde y|_{S_T}$, where $\tilde y$ is the unique extension of $y\in\mathcal Y^\an$ to $S_Y^\nu$
(Lemma~\ref{lem:gl-model-change}), is a continuous surjection $\mathcal Y^\an\to\mathcal T^\an$ with $h\circ\beta$ equal to the restriction to $S_X$; FS
functions of class $a$ on $\mathcal T$ pull back to FS functions of class $(a,0)$ on $\mathcal Y$.
\item There are $e\ge1$, a projectively normal $T'=V(I_{T'})\subseteq\PP^{N_T}$ presenting the Veronese $(S_T)_{(e)}$, forms $G_0,\dots,G_{N'}$ of one
degree on $T'$ with $\bar y_m/\bar y_0=G_m/G_0$ in $k(Y)$, and for the blowup $B$ of $T'$ along $(G_m)$ (Definition~\ref{def:pg-blowup}) homeomorphisms
$\tau\colon\mathcal T^\an\to T'^\an$ and $\vartheta\colon\mathcal Y^\an\to B^\an$ with $\mu\circ\vartheta=\tau\circ\beta$, $\PSH(\mathcal T,et)=\PSH(T',t)\circ\tau$ and
$\PSH(\mathcal Y,(et,s))=\PSH(B,(t,s))\circ\vartheta$ for $t,s\ge0$ (the bridges of Lemma~\ref{lem:gl-bridge-psh} hold for nonnegative classes; FHX uses $t=1/e$, $s=0$).
Projective normality of $X$ is not used in (1)--(3) ($T'$ is projectively normal because $\mathcal T$ is normal); it is used downstream, for finite descent and the bridge for $X$.
\end{enumerate}
\end{lemma}

\begin{proof}
(1) $S_T$ is graded (Lemma~\ref{lem:wu-graded-closure}(1) for $A=k(Y)$), finite over $S_X$ (Theorem~\ref{thm:bj-normalization-finite} for $S_X[\gamma_1,\dots]$,
$\gamma_i$ integral generators of $k(Y)$ over $k(X)$ times powers of $x_0T_1$), integrally closed by construction, integral over the coordinates; its function field
is $k(Y)$, since every $u\in k(Y)$ is $\alpha/\gamma$ with $\alpha$ integral over $A_0=(S_X)_{(x_0)}$ and $0\ne\gamma\in A_0$, and $\alpha(x_0T_1)^M\in S_T$ for $M\gg0$.
(2) Elements of $S_T$ are integral over $S_X\subseteq S_Y$ and lie in $\operatorname{Frac}S_Y=k(Y)(T_1,T_2)$. The chart algebras of $\mathcal Y$ are integrally
closed (smoothness, Definition~\ref{def:gl-smooth}), so Lemma~\ref{lem:gl-model-change} gives the unique extension; $\beta(y)$ is graded (elements of $S_T$ have
degree $(\ast,0)$) and normalized (the coordinates $x$ are those of $\mathcal Y$). Surjectivity follows from (3): $\beta=\tau^{-1}\circ\mu\circ\vartheta$ and $\mu$ is
surjective (Lemma~\ref{lem:pg-blowup-surjective}). Pullback of FS functions: a section of degree $a$ of $\mathcal T$ is an element of $S_Y^\nu$ of degree $(a,0)$, a section of $\mathcal Y$ by Lemma~\ref{lem:gl-model-change}. (3)
$(S_T)_0=k$, so Lemma~\ref{lem:wu-veronese} gives $e$ with $(S_T)_{(e)}$ generated in degree $1$ and integrally closed; present it as $T'$ (Remark~\ref{rmk:gl-bridges}),
giving $\tau$ (Lemmas~\ref{lem:gl-veronese}, \ref{lem:gl-model-change}). Since $k(T)=k(Y)$, the ratios $\bar y_m/\bar y_0$ are $G_m/G_0$ for elements $G_m$ of one degree
$d$ of $(S_T)_{(e)}$ (common denominator), lifted to forms. The bigraded ring of $B$ maps to $S_Y^\nu$ by $x'_i\mapsto$ the degree-$e$ generators of $S_T$ and
$y_m\mapsto\bar y_m$; a bihomogeneous $F$ of bidegree $(a,b)$ maps to $(\bar y_0/G_0)^bF(x',G)$, so the kernel is the ideal $I'$ of Definition~\ref{def:pg-blowup}
(Lemma~\ref{lem:pg-blowup-basic}(1)). The image contains the $(e,1)$-Veronese of $S_Y$ and lies in its normalization, so Lemmas~\ref{lem:gl-veronese} (with the degree
rescaling $(e,1)$, proved in the same way) and~\ref{lem:gl-model-change} give $\vartheta$ and the correspondence of psh classes (Lemma~\ref{lem:gl-model-change} needs the sections
of the larger model to be sections of $\mathcal Y_{(e,1)}$; this follows from the ideal argument on the smooth, hence normal, $\mathcal Y$, without checking that the charts of
$\mathcal Y_{(e,1)}$ are integrally closed); $\mu\circ\vartheta=\tau\circ\beta$ holds on
generators.
\end{proof}

\begin{theorem}[Fibre constancy]\label{thm:gl-fibre-constancy}
\lean{EnvelopeProperty.blowupFibreConstancy_proved, EnvelopeProperty.ProjSubvariety.BlowupCenter.fc_const_on_fibre, EnvelopeProperty.ProjSubvariety.BlowupCenter.IsPSH.apply_eq_of_proj_eq, EnvelopeProperty.ProjSubvariety.BlowupCenter.IsPSH.factorsThrough_proj, EnvelopeProperty.ProjSubvariety.BlowupCenter.IsPSH.exists_eq_comp_proj}
\leanok
\uses{def:pg-blowup, def:pg-blowup-psh, lem:pg-blowup-basic,thm:fc-zariski,thm:fc-analytic-constancy-completed-geometrically-connected}
\source{fhx-envelope-20261028:thm:normal-ample-EP, fhx-envelope-20261028:lem:local-basic, fhx-envelope-20261028:lem:global-to-local, gr25-classically-psh:thm:maximum.principle}
Let $T'\subseteq\PP^{N_T}$ be projectively normal, $B$ a blowup of $T'$ as in Definition~\ref{def:pg-blowup}, $a\in\Q_{>0}$, and
$\Phi\in\PSH(B,(a,0))$. Then $\Phi$ is constant on every fibre of $\mu\colon B^\an\to T'^\an$.
\end{theorem}

\begin{remark}\label{rmk:gl-fibre-constancy}
\uses{thm:gl-fibre-constancy, cor:normal-weak-global}
This is the claim ``$\Phi$ is constant on every fibre of $\beta^\an$'' in Step~1 of the proof of \cite[Theorem~6.1]{FHX}. Its proof there: the fibre over $t$ is
the analytification of the proper $\mathcal H(t)$-scheme $Y\times_T\Spec\mathcal H(t)$, connected because $\beta$ has geometrically connected fibres, hence
compact, connected and without boundary (\cite[\S\S3.4--3.5]{Berkovich90}); $\Phi-m^{-1}\log|\beta^*s|$ is locally psh (\cite[Lemma~5.4(2)]{FHX}) and
its restriction to the fibre is classically psh (\cite[Lemma~5.2(2),(4)]{FHX}); $m^{-1}\log|\beta^*s|$ is constant on the fibre; the maximum set is closed by
upper semicontinuity and open by the local maximum principle \cite[Theorem~6.1]{GR25} (\texttt{thm:maximum.principle}: a classically psh function with a
local maximum at an interior point is locally constant there). \cite{FHX} applies it to the birational part $\beta$ of the Stein factorization of an alteration;
the statement above, for every blowup of a projectively normal variety, is justified by the same argument, as $\mu_*\mathcal O=\mathcal O_{T'}$ (Zariski's main
theorem) gives geometrically connected fibres. The checked native proof now derives every
geometric input from actual projective formal functions and the original whole coefficient
map, preserving the nonreduced closed fibre. Genuine decreasing finite-Fubini--Study nets,
whole native fibre comparisons and the proper maximum principle derive the analytic
constancy without countability. The original compatibility class is provided by an actual
instance and is no longer an unproved mathematical input. The intrinsic BJ/model and
sequence/classical comparisons remain separate obligations.
\end{remark}

\begin{proposition}[Ample envelopes over countable fields; Step~1]\label{prop:gl-ample-countable}
\lean{EnvelopeProperty.ProjSubvariety.Alteration.hasEnvelopeProperty_steinModel, EnvelopeProperty.ProjSubvariety.hasEnvelopeProperty_one_of_isProjNormal_of_countable, EnvelopeProperty.ProjSubvariety.hasEnvelopeProperty_one_of_isProjNormal_of_subfield}
\leanok
\uses{thm:gl-dejong, lem:gl-stein, thm:gl-fibre-constancy, cor:gl-smooth-ep-nef, thm:pg-psh-descent, thm:gl-finite-descent, lem:gl-pullback, lem:gl-ep-basic}
\source{fhx-envelope-20261028:thm:normal-ample-EP}
Let $k$ be countable, and assume the hypotheses of Theorem~\ref{thm:pg-psh-descent} (the imports of Theorem~\ref{thm:normal-local-regularization}) and
Theorems~\ref{thm:gl-dejong}, \ref{thm:gl-fibre-constancy}, \ref{thm:gl-smooth-ep}. If $X\subseteq\PP^N$ is projectively normal, then $\mathrm{EP}_X(1)$.
\end{proposition}

\begin{proof}
Take $\mathcal Y,\mathcal T,h,\beta,T',B,\tau,\vartheta,e$ from Theorem~\ref{thm:gl-dejong} and Lemma~\ref{lem:gl-stein}. We show $\mathrm{EP}_{\mathcal T}(1)$; then
$\mathrm{EP}_X(1)$ by Theorem~\ref{thm:gl-finite-descent}. Let $(\varphi_i)$ be a nonempty family in $\PSH(\mathcal T,1)$ bounded above and $g=\sup_i\varphi_i$. By
Lemma~\ref{lem:gl-stein}(2), $\varphi_i\circ\beta\in\PSH(\mathcal Y,(1,0))$, so $\Phi\coloneqq(g\circ\beta)^\star\in\PSH(\mathcal Y,(1,0))$ by
Corollary~\ref{cor:gl-smooth-ep-nef}. By Lemma~\ref{lem:gl-stein}(3), $\Phi\circ\vartheta^{-1}\in\PSH(B,(1/e,0))$, which is constant on the fibres of $\mu$
(Theorem~\ref{thm:gl-fibre-constancy}); as $\mu\circ\vartheta=\tau\circ\beta$, $\Phi$ is constant on the fibres of $\beta$, and $\beta$ is surjective, so $\Phi=\varphi\circ\beta$
for a function $\varphi$ on $\mathcal T^\an$. Then $(\varphi\circ\tau^{-1})\circ\mu=\Phi\circ\vartheta^{-1}\in\PSH(B,(1/e,0))$, and Theorem~\ref{thm:pg-psh-descent} (with
$a=1/e$) gives $\varphi\circ\tau^{-1}\in\PSH(T',1/e)$, i.e.\ $\varphi\in\PSH(\mathcal T,1)$. Finally $\varphi=g^\star$: from $\Phi\ge g\circ\beta$, $\varphi\ge g$, hence $\varphi\ge
g^\star$ ($\varphi$ is usc); from $\Phi\le g^\star\circ\beta$ (usc and $\ge g\circ\beta$) and surjectivity, $\varphi\le g^\star$.
\end{proof}

\begin{lemma}[Descent to a countable field]\label{lem:gl-field-descent}
\lean{EnvelopeProperty.ProjSubvariety.exists_descendField, EnvelopeProperty.ProjSubvariety.IsProjNormal.descendField, EnvelopeProperty.ProjSubvariety.exists_restrictField_eq_of_isDivisorial, EnvelopeProperty.ProjSubvariety.descendField, EnvelopeProperty.ProjSubvariety.restrictField, EnvelopeProperty.ProjSubvariety.IsPSH.comp_restrictField, EnvelopeProperty.exists_isAlgClosed_countable_subfield}
\leanok
\uses{def:gl-maps, lem:c0-base-change-dvr, cor:c0-base-change-orders, lem:c0-tensor-domain, def:gl-div}
\source{fhx-envelope-20261028:thm:normal-ample-EP}
Let $X=V(I)\subseteq\PP^N_k$ and let $\Phi$ be a countable set of forms. There is a countable algebraically closed $k_0\subseteq k$ such that
$I_0=I\cap k_0[x]$ generates $I$, $\Phi\subseteq k_0[x]$, and $S=S_0\otimes_{k_0}k$. If $S$ is integrally closed, so is $S_0$. For every
$v_0\in X_0^{\divv}$ and $k_0$-linearly independent $c_1,\dots,c_n\in k$ there is $v\in X^\an$ with $r(v)=v_0$ and
$|\sum_\beta c_\beta p_\beta|_v=\max_\beta|p_\beta|_{v_0}$ for all $p_\beta\in S_0$.
\end{lemma}

\begin{proof}
Let $k_0$ be the algebraic closure in $k$ of the field generated by the coefficients of finitely many generators of $I$ and of $\Phi$; it is countable. $I_0k[x]=I$
and $I_0$ is prime ($k[x]$ is faithfully flat over $k_0[x]$). Normality descends: for $z=\alpha/\gamma\in\operatorname{Frac}S_0$ integral over $S_0$, $z\in S$, i.e.\
$\alpha\in\gamma S\cap S_0=\gamma S_0$ by faithful flatness. Let $R'_0\supseteq S_0[t]$ and $c$ be the data of $v_0$ (Definition~\ref{def:gl-div}) and
$R'=(R'_0\otimes_{k_0}k)_{\mathfrak m\otimes k}$ (Lemma~\ref{lem:c0-base-change-dvr}), a discrete valuation ring containing $S[t]$ with
$\operatorname{ord}_{R'}(\sum c_\beta\otimes g_\beta)=\min_\beta\operatorname{ord}_{R'_0}(g_\beta)$ (Corollary~\ref{cor:c0-base-change-orders}). Put
$|p|_v=e^{-c\operatorname{ord}_{R'}(p)}$ for $p\in S$; it is multiplicative, trivial on $k$, graded (components are computed coefficientwise in a $k_0$-basis),
normalized ($x_i\in S_0$), restricts to $v_0$, and has the displayed values.
\end{proof}

\begin{proposition}[Ample envelopes; Step~2]\label{prop:gl-ample-rational}
\lean{EnvelopeProperty.ProjSubvariety.hasEnvelopeProperty_one_of_countable, EnvelopeProperty.ProjSubvariety.hasEnvelopeProperty_one_of_forall_countable, EnvelopeProperty.ProjSubvariety.envelope_comp_restrictField}
\leanok
\uses{prop:gl-ample-countable, lem:gl-field-descent, prop:gl-pl-envelope, lem:gl-fs-envelope, thm:gl-comparison, lem:gl-pullback}
\source{fhx-envelope-20261028:thm:normal-ample-EP}
Under the hypotheses of Proposition~\ref{prop:gl-ample-countable} for every countable algebraically closed subfield (and $k$ arbitrary), if $X\subseteq\PP^N$ is
projectively normal then $\mathrm{EP}_X(1)$.
\end{proposition}

\begin{proof}
By Proposition~\ref{prop:gl-pl-envelope} it suffices that $\mathrm P_1(\phi)$ is psh for $\phi\in\mathrm{PL}(X)$. Choose $k_0$ (Lemma~\ref{lem:gl-field-descent}) for the
forms defining $\phi$, so that $\phi=\phi_0\circ r$ with $\phi_0$ continuous on $X_0^\an$: write $\phi$ as a difference of FS functions $\phi_1-\phi_2$, bounded
below by $-C_1$, $-C_2$, and put $\phi_0=\max(\phi_{1,0},-C_1)-\max(\phi_{2,0},-C_2)$ with the descended FS functions $\phi_{i,0}$ (it is not
needed that $\phi_0$ is PL, i.e.\ that base point freeness descends); $X_0$ is projectively normal and $\mathrm{EP}_{X_0}(1)$ holds by
Proposition~\ref{prop:gl-ample-countable}, so $\mathrm P_1(\phi_0)\in\PSH(X_0,1)$ (envelopes of continuous obstacles), and $\mathrm P_1(\phi_0)\circ r\in\PSH(X,1)$
(Lemma~\ref{lem:gl-pullback}) is $\le\phi$, hence $\le\mathrm P_1(\phi)$. Conversely let $\psi=m^{-1}\max_t\{\log|s_t|+\lambda_t\}\in\mathcal H_{\mathrm{fin}}(X,1)$ with
$\psi\le\phi$; write $s_t=\sum_\beta c_\beta s_{t\beta}$ with $s_{t\beta}\in(S_0)_m$ and $c_\beta\in k$ linearly independent over $k_0$, and
$\psi_0=m^{-1}\max_{t,\beta}\{\log|s_{t\beta}|+\lambda_t\}\in\mathcal H(X_0,1)$. Then $\psi\le\psi_0\circ r$ (ultrametric inequality). For $v_0\in X_0^{\divv}$ the point $v$ of
Lemma~\ref{lem:gl-field-descent} gives $\psi_0(v_0)=\psi(v)\le\phi(v)=\phi_0(v_0)$; by Theorem~\ref{thm:gl-comparison} on $X_0$, $\psi_0\le\phi_0$, so
$\psi_0\le\mathrm P_1(\phi_0)$ and $\psi\le\mathrm P_1(\phi_0)\circ r$. By Lemma~\ref{lem:gl-fs-envelope}, $\mathrm P_1(\phi)=\mathrm P_1(\phi_0)\circ r$, which is psh.
\end{proof}

\begin{theorem}[Ample classes on normal varieties]\label{thm:gl-normal-ample-ep}
\lean{EnvelopeProperty.MultiModel.hasEnvelopeProperty_of_isNormal_of_pos, EnvelopeProperty.MultiModel.hasEnvelopeProperty_of_isNormal_of_pos_of_forall_countable, EnvelopeProperty.MultiModel.hasEnvelopeProperty_one_of_isNormal_of_forall_countable, EnvelopeProperty.MultiModel.hasEnvelopeProperty_of_rank_zero, EnvelopeProperty.MultiModel.eq_triv_of_rank_zero, EnvelopeProperty.MultiModel.hasEnvelopeProperty_of_isNormal_of_pos_of_forall_subfield}
\leanok
\uses{prop:gl-ample-rational, lem:gl-veronese, rmk:gl-bridges, lem:gl-ep-basic, lem:wu-veronese}
\source{fhx-envelope-20261028:thm:normal-ample-EP}
Under the hypotheses of Proposition~\ref{prop:gl-ample-rational}, $\mathrm{EP}_{\mathcal X}(\theta)$ holds for every normal model $\mathcal X$ and every positive
$\theta\in\R^r$.
\end{theorem}

\begin{proof}
Step~3 of \cite{FHX}: by Lemma~\ref{lem:gl-ep-basic}(1) with rational positive $\theta_j\downarrow\theta$, and (2), it suffices to treat positive $a\in\N^r$. The model
$\mathcal X_{(a)}$ is normal of rank~$1$ with degree-zero part $k$ (elements of $S_0$ are algebraic over $k$, and $k$ is algebraically closed in $K$); a
Veronese of it is generated in degree $1$ and integrally closed (Lemma~\ref{lem:wu-veronese}), i.e.\ a projectively normal $V(I)$
(Remark~\ref{rmk:gl-bridges}). Proposition~\ref{prop:gl-ample-rational} and Lemma~\ref{lem:gl-veronese} conclude: $\mathrm{EP}_{V(I)}(1)$ gives
$\mathrm{EP}$ of the Veronese $(\mathcal X_{(a)})_{(e)}$ at $1$, hence $\mathrm{EP}_{\mathcal X_{(a)}}(e)$, hence $\mathrm{EP}_{\mathcal X_{(a)}}(1)$ by
Lemma~\ref{lem:gl-ep-basic}(2). For $r=0$ there is nothing to prove: $S=S_0=k$, so $\mathcal X^\an=\{\star\}$ and every bounded family has a constant, hence psh,
regularized supremum. The ambient dimension $N$ of $V(I)$ depends on $a$, so the hypothesis of Proposition~\ref{prop:gl-ample-rational} is used for all $N$.
\end{proof}

\subsection{Homogeneous division and arbitrary classes}\label{subsec:gl-siu}

\begin{lemma}[Points attached to prime divisors]\label{lem:gl-divisor-points}
\lean{EnvelopeProperty.MultiModel.primeDivPoint, EnvelopeProperty.MultiModel.divisorVal, EnvelopeProperty.MultiModel.primeOrd, EnvelopeProperty.MultiModel.secAbs_primeDivPoint, EnvelopeProperty.MultiModel.divisorVal_laurentMono_le_one, EnvelopeProperty.MultiModel.primeDivPoint_mem_valSet, EnvelopeProperty.MultiModel.hasGenericSupport_primeDivPoint, EnvelopeProperty.MultiModel.isDVR_localization_of_height_eq_one, EnvelopeProperty.MultiModel.isDVR_localization_iff_height_eq_one, EnvelopeProperty.laurentMono}
\leanok
\uses{def:gl-points, lem:gl-multigraded-extension, lem:gl-sections, lem:bj-dvr-minimal, lem:pg-div-points}
\source{fhx-envelope-20261028:lem:fixed-coefficient-division, bj22-trivval:lem:FShomdiv}
Let $\mathcal X$ be normal and $P$ a height-one multihomogeneous prime of $S$ that is \emph{relevant} (no block of coordinates lies in $P$). Then $S_P$ is a
discrete valuation ring with valuation $\operatorname{ord}_P$ on $K(T)$; $\operatorname{ord}_P(fT^a)\ge0$ for $f\in R_a$; and $|gT^a|_{v_P}=\exp(-\operatorname{ord}_P(gT^a))$,
$g\in S_a$, defines $v_P\in\mathcal X^{\mathrm{val}}$ with $|f|_{v_P}=\exp(-\operatorname{ord}_P(fT^a))$ for $f\in R_a$.
\end{lemma}

\begin{proof}
$P$ is minimal over each of its nonzero elements, so $S_P$ is a discrete valuation ring (Lemma~\ref{lem:bj-dvr-minimal}). In each block a coordinate
$z^{(c)}_{i_c}\notin P$ has order $0$ and all coordinates have order $\ge0$, so the hypotheses of Lemma~\ref{lem:gl-multigraded-extension} hold (as in
Lemma~\ref{lem:pg-div-points}). For $f\in R_a$ and the monomial $z^\beta$ of degree $b$ in the $z^{(c)}_{i_c}$, $\operatorname{ord}_P(fT^a)=
\operatorname{ord}_P(f\tilde z^\beta T^{a+b})\ge0$, and $|f|_{v_P}=\max_\beta\exp(-\operatorname{ord}_P(fT^a)-\operatorname{ord}_P(z^\beta))=\exp(-\operatorname{ord}_P(fT^a))$.
\end{proof}

\begin{lemma}[Sections on normal models]\label{lem:gl-section-criterion}
\lean{EnvelopeProperty.MultiModel.mem_sections_of_forall_divisorVal_le_one, EnvelopeProperty.MultiModel.mem_sections_of_forall_height_eq_one, EnvelopeProperty.MultiModel.isMultihomogeneousIdeal_of_mem_minimalPrimes, EnvelopeProperty.MultiModel.isPrime_homCore}
\leanok
\uses{lem:gl-divisor-points, lem:bj-dvr-minimal, def:gl-sections}
\source{bj22-trivval:lem:FShomdiv, fhx-envelope-20261028:lem:fixed-coefficient-division}
Let $\mathcal X$ be normal, $a\in\mathbf Z^r$ and $f\in K$. If $\operatorname{ord}_P(fT^a)\ge0$ for every relevant height-one multihomogeneous prime $P$, then
$f\in R_a$.
\end{lemma}

\begin{proof}
Only finitely many height-one primes $Q$ of the Noetherian normal domain $S$ have $\operatorname{ord}_Q(fT^a)<0$. Choose $b\in\N^r$ with $a+b\ge0$ and
$b_c\ge-\operatorname{ord}_Q(fT^a)$ for every such $Q$ containing all coordinates of the block $c$. For $\beta\in\mathrm{Mon}(b)$, $p=f\tilde z^\beta T^{a+b}$ is
homogeneous in $\operatorname{Frac}S$, so $p=A/B$ with $A,B\in S$ homogeneous (take the extremal homogeneous component of a denominator). The minimal primes $Q$ of
$BS$ are homogeneous of height one, and $\operatorname{ord}_Q(p)=\operatorname{ord}_Q(fT^a)+\operatorname{ord}_Q(z^\beta)\ge0$: for relevant $Q$ by hypothesis, for
irrelevant $Q$ (containing a whole block $c$) because $\operatorname{ord}_Q(z^\beta)\ge b_c$. Lemma~\ref{lem:bj-dvr-minimal} gives $A\in BS$, i.e.\ $p\in S$.
\end{proof}

\begin{lemma}[Division with a fixed coefficient]\label{lem:gl-fixed-coefficient-division}
\lean{EnvelopeProperty.MultiModel.IsPSHHom.exists_eq_add_of_le_logSecAbs, EnvelopeProperty.MultiModel.FSData.fixedDiv, EnvelopeProperty.MultiModel.FSData.eval_fixedDiv}
\leanok
\uses{lem:gl-divisor-points, lem:gl-section-criterion, cor:gl-pshhom, thm:gl-comparison, lem:gl-div-generic, lem:gl-sections}
\source{fhx-envelope-20261028:lem:fixed-coefficient-division}
Let $\mathcal X$ be normal, $\eta\in\R^r$, $e\in\Q^r$ and $0\ne s_E\in R_{me}$ with $me\in\mathbf Z^r$; put $e_E=m^{-1}\log|s_E|$. If
$\psi\in\PSH_{\mathrm{hom}}(\mathcal X,\eta)$ and $\psi\le e_E$, then for every rational $0<t<1$ there is $q_t\in\PSH_{\mathrm{hom}}(\mathcal X,\eta-te)$ with
$\psi=te_E+q_t$.
\end{lemma}

\begin{proof}
Let $P_1,\dots,P_s$ be the relevant height-one multihomogeneous primes with $\alpha_\ell\coloneqq m^{-1}\operatorname{ord}_{P_\ell}(s_ET^{me})>0$; at all other relevant ones
the order is $0$ (Lemma~\ref{lem:gl-divisor-points}). By Corollary~\ref{cor:gl-pshhom}, $\psi$ is the limit of a decreasing net $\psi_\alpha\in\mathcal
H_{\mathrm{hom}}(a_\alpha)$, $a_\alpha\to\eta$. Since $\psi_\alpha(v_{P_\ell})\downarrow\psi(v_{P_\ell})\le e_E(v_{P_\ell})=-\alpha_\ell<-t\alpha_\ell$, for $\alpha$ large
$\psi_\alpha(v_{P_\ell})\le-t\alpha_\ell$ for all $\ell$. Write $\psi_\alpha=M^{-1}\max_i\log|\sigma_i|$ with $\sigma_i\in R_{Ma_\alpha}$ and $Mt/m\in\mathbf Z$. Then
$\operatorname{ord}_{P_\ell}(\sigma_iT^{Ma_\alpha})\ge Mt\alpha_\ell=(Mt/m)\operatorname{ord}_{P_\ell}(s_ET^{me})$, so $\sigma_i/s_E^{Mt/m}$ has nonnegative order at every relevant
$P$, and lies in $R_{Ma_\alpha-Mte}$ by Lemma~\ref{lem:gl-section-criterion}. By multiplicativity (Lemma~\ref{lem:gl-sections}(2)) $\psi_\alpha=te_E+q_\alpha$ on
$\mathcal X^\an$ with $q_\alpha=M^{-1}\max_i\log|\sigma_i/s_E^{Mt/m}|\in\mathcal H_{\mathrm{hom}}(a_\alpha-te)$. On $\Omega=\{e_E>-\infty\}\supseteq\mathcal X^{\divv}$
(Lemma~\ref{lem:gl-div-generic}) the $q_\alpha=\psi_\alpha-te_E$ decrease, hence they decrease everywhere (Theorem~\ref{thm:gl-comparison}); $q_\alpha(\star)=0$, so
$q_t=\lim q_\alpha\in\PSH_{\mathrm{hom}}(\eta-te)$ by definition (no PSH2 is needed: the $q_\alpha$ are Fubini--Study functions). In the limit $\psi=te_E+q_t$ on
$\Omega$, and off $\Omega$ both sides are $-\infty$ ($\psi\le e_E$).
\end{proof}

\begin{theorem}[Homogeneous Siu division]\label{thm:gl-homogeneous-siu}
\lean{EnvelopeProperty.MultiModel.IsPSHHom.exists_isPSHHom_eq_add_of_le_logSecAbs}
\leanok
\uses{lem:gl-fixed-coefficient-division, thm:gl-normal-ample-ep, lem:gl-usc-local, thm:gl-comparison, lem:gl-scaling, thm:gl-fs-system-theory, lem:gl-div-generic}
\source{fhx-envelope-20261028:thm:homogeneous-Siu, fhx-envelope-20261028:rem:BJ-6-21}
Assume the hypotheses of Theorem~\ref{thm:gl-normal-ample-ep}. Let $\mathcal X$ be normal, $\theta\in N(\mathcal X)$, and $s_E$, $e$, $e_E$ as in
Lemma~\ref{lem:gl-fixed-coefficient-division}. If $\psi\in\PSH_{\mathrm{hom}}(\mathcal X,\theta+e)$ and $\psi\le e_E$, then $\psi=e_E+\varphi$ with
$\varphi\in\PSH_{\mathrm{hom}}(\mathcal X,\theta)$.
\end{theorem}

\begin{proof}
Let $\Omega=\{e_E>-\infty\}$ (open) and $u=\psi-e_E$ on $\Omega$, usc and $\le0$ there. Choose $N\in\N$, $N\ge1$, with $\theta+N\mathbf1$ positive, put $A=N\mathbf 1$, and for
$\beta\in\mathrm{Mon}(A)$ let $e_\beta=\log|\tilde z^\beta|\in\mathcal H_{\mathrm{hom}}(A)$ and $\Omega_\beta=\{e_\beta>-\infty\}$; $\max_\beta e_\beta=0$. Let
$V=\Omega\cap\bigcap_\beta\Omega_\beta\supseteq\mathcal X^{\divv}$ (Lemma~\ref{lem:gl-div-generic}). Fix $\beta$ and rational $0<\epsilon<1$, and choose $t_0<1$ with
$(1-t_0)e\le\epsilon A$.

For rational $t\in[t_0,1)$, Lemma~\ref{lem:gl-fixed-coefficient-division} applied to $\psi+e_\beta\in\PSH_{\mathrm{hom}}(\theta+e+A)$ (which is $\le e_E$) gives
$q_t\in\PSH_{\mathrm{hom}}(\theta+A+(1-t)e)\subseteq\PSH_{\mathrm{hom}}(\theta+(1+\epsilon)A)$ with $\psi+e_\beta=te_E+q_t$, i.e.\ $q_t=u+e_\beta+(1-t)e_E$ on $\Omega$. On
$\Omega$, $q_t\le e_\beta$; by Lemma~\ref{lem:gl-usc-local}(3), $q_t\le e_\beta$ everywhere. Let $g=\sup_tq_t\le e_\beta$. By Theorem~\ref{thm:gl-normal-ample-ep} at the
positive class $\theta+(1+\epsilon)A$, $\chi_\beta\coloneqq g^\star\in\PSH(\theta+(1+\epsilon)A)$; it is homogeneous (Lemma~\ref{lem:gl-scaling}(5)) and $\le e_\beta$ ($e_\beta$ is
continuous). On $\Omega$, $g=u+e_\beta$ (let $t\uparrow1$), which is usc on $\Omega$, so $\chi_\beta=u+e_\beta$ on $\Omega$ (Lemma~\ref{lem:gl-usc-local}(1)).

Lemma~\ref{lem:gl-fixed-coefficient-division} applied to $\chi_\beta$, the section $\tilde z^\beta$ of degree $A$ and $t=1-\epsilon$ gives
$q'_\beta\in\PSH_{\mathrm{hom}}(\theta+2\epsilon A)$ with $\chi_\beta=(1-\epsilon)e_\beta+q'_\beta$, so $q'_\beta=u+\epsilon e_\beta$ on $V$. Then
$\varphi_\epsilon=\max_\beta q'_\beta\in\PSH_{\mathrm{hom}}(\theta+2\epsilon A)$ equals $u+\epsilon\max_\beta e_\beta=u$ on $V$. By Lemma~\ref{lem:gl-usc-local}(3) the
$\varphi_\epsilon$ do not depend on $\epsilon$; call them $\varphi$. PSH2 at $\theta\in N(\mathcal X)$ (Theorem~\ref{thm:gl-fs-system-theory}(2), constant sequence,
classes $\theta+2\epsilon A\to\theta$) gives $\varphi\in\PSH_{\mathrm{hom}}(\theta)$. Finally $e_E+\varphi$ and $\psi$ are $(\theta+e)$-psh and agree on $V$, hence everywhere.
\end{proof}

\begin{proposition}[Envelopes from homogeneous slices]\label{prop:gl-homogeneous-slices}
\lean{EnvelopeProperty.MultiModel.hasEnvelopeProperty_of_isPSHHom_uscReg_iSup}
\leanok
\uses{thm:gl-homog-decomposition, thm:gl-fs-system-theory, lem:gl-scaling}
\source{fhx-envelope-20261028:prop:homogeneous-slices, fhx-envelope-20261028:eq:Legendre-reconstruction}
Let $\theta\in N(\mathcal X)$. Suppose that $(\sup_i\psi_i)^\star\in\PSH_{\mathrm{hom}}(\mathcal X,\theta)$ for every nonempty increasing net $(\psi_i)$ in
$\PSH_{\mathrm{hom}}(\mathcal X,\theta)$. Then $\mathrm{EP}_{\mathcal X}(\theta)$.
\end{proposition}

\begin{proof}
Let $(\varphi_i)$ be a nonempty family in $\PSH(\theta)$ bounded above; replacing it by its finite maxima we may assume it is an increasing net, and after subtracting a
constant that $b_i=\varphi_i(\star)$ increases to $0$. Let $f=\sup_i\varphi_i$. For $\lambda<0$ let $I_\lambda=\{i:b_i>\lambda\}$ (cofinal), $\psi_{i,\lambda}=(\hat\varphi_i)_\lambda
\in\PSH_{\mathrm{hom}}(\theta)$ for $i\in I_\lambda$ (Theorem~\ref{thm:gl-homog-decomposition}), increasing in $i$, $g_\lambda=\sup_{i\in I_\lambda}\psi_{i,\lambda}$ and
$\psi_\lambda=g_\lambda^\star\in\PSH_{\mathrm{hom}}(\theta)$ (hypothesis). The $\psi_\lambda$ decrease in $\lambda$ (slices decrease in $\lambda$, $I_{\lambda'}\subseteq I_\lambda$ for
$\lambda\le\lambda'$). Let $\Psi=(\sup_n\psi_{-n})^\star\in\PSH_{\mathrm{hom}}(\theta)$ (hypothesis); $\Psi\ge\psi_\lambda$ for all $\lambda<0$. By
Theorem~\ref{thm:gl-homog-decomposition}, $f=\sup_i\sup_{\lambda<b_i}(\psi_{i,\lambda}+\lambda)=\sup_{\lambda<0}(g_\lambda+\lambda)$ pointwise.

For $n\ge1$ let $h_n=2^{-n}$, $A_n=\{-n+\ell h_n:0\le\ell<n2^n\}$ and $U_n=\max\{\Psi-n,\max_{a\in A_n}(\psi_a+a+h_n)\}\in\PSH(\theta)$, with $U_n(\star)=0$. As in
\cite{FHX}, $U_{n+1}\le U_n$ everywhere: for $a'\in A_{n+1}$ in $[a,a+h_n)$, $a\in A_n$, $\psi_{a'}+a'+h_{n+1}\le\psi_a+a+h_n$; for $a'\in[-n-1,-n)$,
$\psi_{a'}+a'+h_{n+1}\le\Psi-n$. By PSH2 at $\theta$, $U=\lim U_n\in\PSH(\theta)$. Lower bound: for $\lambda<-n$, $\psi_\lambda+\lambda\le\Psi-n$, and for $\lambda\in[a,a+h_n)$,
$\psi_\lambda+\lambda\le\psi_a+a+h_n$; so $f\le\sup_\lambda(\psi_\lambda+\lambda)\le U_n$, and $f^\star\le U$. Upper bound: $\psi_\lambda+\lambda=(g_\lambda+\lambda)^\star\le f^\star$, so
$U_n\le\max\{\Psi-n,f^\star+h_n\}$, and $U\le f^\star$ ($\Psi<\infty$). Hence $f^\star=U\in\PSH(\theta)$.
\end{proof}

\begin{corollary}[Arbitrary nef classes on normal varieties]\label{cor:gl-normal-ep}
\lean{EnvelopeProperty.MultiModel.hasEnvelopeProperty_of_isNormal}
\leanok
\uses{prop:gl-homogeneous-slices, thm:gl-homogeneous-siu, thm:gl-normal-ample-ep, lem:gl-usc-local, lem:gl-scaling}
\source{fhx-envelope-20261028:cor:normal-EP}
Under the hypotheses of Theorem~\ref{thm:gl-normal-ample-ep}, $\mathrm{EP}_{\mathcal X}(\theta)$ holds for every normal model $\mathcal X$ and every $\theta\in N(\mathcal X)$.
\end{corollary}

\begin{proof}
We verify the hypothesis of Proposition~\ref{prop:gl-homogeneous-slices}. Let $(\psi_i)$ be an increasing net in $\PSH_{\mathrm{hom}}(\theta)$. Choose $N$ with
$\eta=\theta+N\mathbf1$ positive, a monomial $s_E=\tilde z^{\beta_0}$ of degree $e=N\mathbf 1$, $e_E=\log|s_E|$, $\Omega=\{e_E>-\infty\}$. The $\psi_i+e_E\in
\PSH_{\mathrm{hom}}(\eta)$ increase and are $\le e_E\le0$; by Theorem~\ref{thm:gl-normal-ample-ep}, $\chi=(\sup_i(\psi_i+e_E))^\star\in\PSH_{\mathrm{hom}}(\eta)$, and
$\chi\le e_E$. Theorem~\ref{thm:gl-homogeneous-siu} gives $\chi=e_E+\psi$ with $\psi\in\PSH_{\mathrm{hom}}(\theta)$. On $\Omega$, $\chi=(\sup_i\psi_i)^\star+e_E$
(Lemma~\ref{lem:gl-usc-local}(2)), so $\psi=G\coloneqq(\sup_i\psi_i)^\star$ on $\Omega\supseteq\mathcal X^{\divv}$. Then $\psi_i\le\psi$ on $\Omega$, hence everywhere, so
$G\le\psi$; and $\psi\le G$ everywhere since $G$ is usc and they agree on $\Omega$ (Lemma~\ref{lem:gl-usc-local}(3)). So $G=\psi$.
\end{proof}

\begin{remark}[Deviations in \S7]\label{rmk:gl-siu-source}
\uses{thm:gl-homogeneous-siu, prop:gl-homogeneous-slices, cor:gl-normal-ep}
\cite{FHX} uses ``divisorial values'' and the direction envelope property $\Rightarrow$ convergence on $X^{\divv}$ of \cite[Proposition~2.4]{FHX}, which rests on
\cite[Theorem~5.6]{BJ22}, and the finiteness of psh functions on $X^{\divv}$ (\cite[Corollary~4.17]{BJ22}). Here the increasing limits are taken as usc regularizations,
and divisorial values are replaced by values on the open sets $\Omega$, $V$, where the usc regularization of a supremum of functions that are usc on the open set
does nothing (Lemma~\ref{lem:gl-usc-local}); the comparison theorem transports equalities from $V\supseteq\mathcal X^{\divv}$ to $\mathcal X^\an$. The hypothesis
of Proposition~\ref{prop:gl-homogeneous-slices} is accordingly the usc form. As in \cite[Remark~7.3]{FHX} (\texttt{rem:BJ-6-21}), the division is only performed with a
fixed coefficient, so \cite[Theorem~6.21]{BJ22} is not used. The restriction $\theta\in N(\mathcal X)$ comes from PSH2 (Remark~\ref{rmk:gl-nonnef}).
\end{remark}

\subsection{Normalization and the main theorem}\label{subsec:gl-main}

\begin{definition}[Unibranch models]\label{def:gl-unibranch}
\lean{EnvelopeProperty.MultiModel.IsUnibranch, EnvelopeProperty.MultiModel.FiniteHom.IsMaxSpecInjective, EnvelopeProperty.MultiModel.isUnibranch_iff, EnvelopeProperty.MultiModel.isUnibranch_iff_anMap_injective, EnvelopeProperty.MultiModel.isUnibranch_tfae, EnvelopeProperty.MultiModel.isUnibranch_of_isNormal}
\leanok
\uses{rmk:gl-model-examples, lem:gl-charts}
\source{fhx-envelope-20261028:lem:unibranch-universal, bj22-trivval:sec:envprop}
Let $\mathcal X$ be a model and $\mathcal X^\nu$ its normalization (Remark~\ref{rmk:gl-model-examples}(2)), with chart algebras $A_j\subseteq B_j$ ($B_j$ finite over
$A_j$, same fraction field). $\mathcal X$ is \emph{unibranch} if for every chart $j$ the map $\operatorname{MaxSpec}B_j\to\operatorname{MaxSpec}A_j$ is injective.
\end{definition}

\begin{lemma}[Radicial extensions]\label{lem:gl-radicial}
\lean{EnvelopeProperty.mem_nilradical_of_forall_algHom_eq_zero, EnvelopeProperty.algHom_eq_of_maxSpec_injective, EnvelopeProperty.ker_lmul'_le_nilradical, EnvelopeProperty.ringHom_ext_of_maxSpec_injective, EnvelopeProperty.injective_comap_of_maxSpec_injective, EnvelopeProperty.bijective_comap_of_maxSpec_injective, EnvelopeProperty.isPurelyInseparable_residueField_of_maxSpec_injective, EnvelopeProperty.isPurelyInseparable_residueField_under, EnvelopeProperty.bijective_algebraMap_residueField_of_maxSpec_injective}
\leanok
\source{fhx-envelope-20261028:lem:unibranch-universal}
Let $A\subseteq B$ be finitely generated $k$-algebras, $B$ finite over $A$, with $\operatorname{MaxSpec}B\to\operatorname{MaxSpec}A$ injective. Then the kernel $J$ of
$B\otimes_AB\to B$ is nil; $\Spec B\to\Spec A$ is bijective; and every residue field extension $\kappa(\mathfrak q)/\kappa(\mathfrak p)$ is purely inseparable (trivial
if $\operatorname{char}k=0$).
\end{lemma}

\begin{proof}
A $k$-point of the finitely generated $k$-algebra $B\otimes_AB$ is a pair of $k$-points of $B$ over one $k$-point of $A$, equal by hypothesis, so it kills every
$b\otimes1-1\otimes b$, which generate $J$. Thus $J$ lies in every maximal ideal, hence in the nilradical ($B\otimes_AB$ is a Jacobson ring; this is the
Nullstellensatz argument of the proof of \cite[Lemma~8.1]{FHX}). Surjectivity of $\Spec B\to\Spec A$ is lying over. If $\mathfrak q,\mathfrak q'$ lie over $\mathfrak p$,
choose a field $L$ and a map $\kappa(\mathfrak q)\otimes_{\kappa(\mathfrak p)}\kappa(\mathfrak q')\to L$; the composite $B\otimes_AB\to L$ kills the nil ideal $J$, so the two maps
$B\to L$ agree, whence $\mathfrak q=\mathfrak q'$, and, taking $\mathfrak q=\mathfrak q'$, any two $\kappa(\mathfrak p)$-embeddings of $\kappa(\mathfrak q)$ into $L$ agree; for the
finite extension $\kappa(\mathfrak q)/\kappa(\mathfrak p)$ this is pure inseparability.
\end{proof}

\begin{lemma}[The normalization of a unibranch model]\label{lem:gl-unibranch-universal}
\lean{EnvelopeProperty.MultiModel.IsUnibranch.anHomeomorph, EnvelopeProperty.MultiModel.FiniteHom.anMap_injective_of_isMaxSpecInjective, EnvelopeProperty.Berk.comap_injective_of_maxSpec_injective, EnvelopeProperty.MultiModel.IsUnibranch.isRadicialIntegral, EnvelopeProperty.MultiModel.IsUnibranch.isSubintegralTower, EnvelopeProperty.MultiModel.IsUnibranch.exists_pow_pow_mem_chartAlg, EnvelopeProperty.MultiModel.IsUnibranch.exists_addPow_mem_chartAlg, EnvelopeProperty.MultiModel.IsUnibranch.isPSH_comp}
\leanok
\uses{def:gl-unibranch, lem:gl-radicial, lem:gl-pullback, lem:gl-charts}
\source{fhx-envelope-20261028:lem:unibranch-universal}
If $\mathcal X$ is unibranch, $\nu^\an\colon(\mathcal X^\nu)^\an\to\mathcal X^\an$ is a homeomorphism, every chart extension $A_j\subseteq B_j$ satisfies the conclusions
of Lemma~\ref{lem:gl-radicial}, and in characteristic $0$ it is subintegral (integral, bijective on spectra, isomorphisms on residue fields).
\end{lemma}

\begin{proof}
$\nu^\an$ is continuous and surjective (Lemma~\ref{lem:gl-pullback}). Let $v\in\mathcal X^\an$ and $j$ a chart with $|z_j|_v=1$. Points over $v$ correspond to points of
$(\Spec B_j)^\an$ over $\iota_j(v)$, i.e.\ to a prime $\mathfrak q$ over the support $\mathfrak p$ of $\iota_j(v)$, unique by Lemma~\ref{lem:gl-radicial}, together with an
absolute value on $\kappa(\mathfrak q)$ extending the one on $\kappa(\mathfrak p)$, unique because $x^{p^n}\in\kappa(\mathfrak p)$ for $x\in\kappa(\mathfrak q)$ (or $\kappa(\mathfrak
q)=\kappa(\mathfrak p)$). A continuous bijection of compact Hausdorff spaces is a homeomorphism.
\end{proof}

\begin{definition}[Radicial integral extensions]\label{def:gl-radicial-integral}
\lean{EnvelopeProperty.Algebra.IsRadicialIntegral, EnvelopeProperty.Algebra.IsRadicialIntegral.of_maxSpec_injective, EnvelopeProperty.Algebra.IsRadicialIntegral.subalgebra, EnvelopeProperty.Algebra.IsRadicialIntegral.over_subalgebra, EnvelopeProperty.Algebra.IsRadicialIntegral.isPurelyInseparable_of_liesOver}
\leanok
\uses{lem:gl-radicial}
An extension $A\to B$ of rings is \emph{radicial integral} if it is injective and integral, $\Spec B\to\Spec A$ is bijective, and every residue field
extension $\kappa(\mathfrak q)/\kappa(\mathfrak q\cap A)$ is purely inseparable. By Lemma~\ref{lem:gl-radicial} this holds for a finite extension $A\subseteq B$ of
finitely generated $k$-algebras with $\operatorname{MaxSpec}B\to\operatorname{MaxSpec}A$ injective. If $A\to B$ is radicial integral and $S$ is an intermediate
$A$-subalgebra, then $A\to S$ and $S\to B$ are radicial integral (lying over for the integral extension $S\subseteq B$, and pure inseparability in towers).
For injective maps this is the ring-level form of a universal homeomorphism (Stacks \texttt{04DF}, \texttt{01S4}); that equivalence is not formalized
and not needed.
\end{definition}

\begin{lemma}[Structure of radicial integral extensions]\label{lem:gl-radicial-structure}
\lean{EnvelopeProperty.exists_sq_cube_of_pow_mem_of_pow_succ_mem, EnvelopeProperty.pow_mem_of_sq_mem_of_cube_mem, EnvelopeProperty.exists_mul_not_mem_and_mul_pow_mem, EnvelopeProperty.exists_not_mem_and_mul_pow_mem_of_localization, EnvelopeProperty.exists_sq_cube_or_prime_of_finite, EnvelopeProperty.Algebra.IsRadicialIntegral.exists_sq_cube_or_prime}
\leanok
\uses{def:gl-radicial-integral}
\source{stacks-universal-homeomorphisms:tag-0CNC}
Let $A\subseteq B$ be radicial integral with $A\ne B$. There is $b\in B\setminus A$ with $b^2,b^3\in A$, or a prime $\ell$ with $\ell b\in A$ and $b^\ell\in A$.
\end{lemma}

\begin{proof}
This is Stacks \texttt{0CNC}, proved at ring level. Pick $b_0\notin A$; replacing $B$ by $A[b_0]$ (Definition~\ref{def:gl-radicial-integral}) makes $B$ finite over
$A$. Let $I=\{a\in A:aB\subseteq A\}$, a proper ideal, and $\mathfrak p$ a minimal prime over $I$. Statements about $A_{\mathfrak p}\subseteq B_{\mathfrak p}$ are read
elementwise ("$ty\in A$ for some $t\notin\mathfrak p$"); this is Stacks \texttt{0CN9}, whose proof Stacks omits. If the unique prime of $B$ over $\mathfrak p$ is not
contained in $A_{\mathfrak p}$, prime avoidance and minimality of $\mathfrak p$ give $f\notin A$ with $f^n,f^{n+1}\in A$, and Stacks \texttt{0CNA} gives $b'\notin A$ with
$b'^2,b'^3\in A$. Otherwise the residue extension at $\mathfrak p$ is nontrivial and purely inseparable of characteristic $\ell$, and a lift of an element outside the
base field with $\ell$-th power inside gives $b$ with $\ell b,b^\ell\in A$.
\end{proof}

\begin{theorem}[Frobenius descent]\label{thm:gl-frobenius-descent}
\lean{EnvelopeProperty.Algebra.IsRadicialIntegral.exists_pow_pow_mem, EnvelopeProperty.Algebra.IsRadicialIntegral.isPRadical, EnvelopeProperty.exists_pow_pow_mem_of_maxSpec_injective}
\leanok
\uses{lem:gl-radicial, def:gl-radicial-integral, lem:gl-radicial-structure}
\source{fhx-envelope-20261028:thm:normalization-reflection, stacks-universal-homeomorphisms:tag-0CNF}
Let $\operatorname{char}k=p>0$ and $A\subseteq B$ a finite extension of finitely generated $k$-algebras such that $\Spec B\to\Spec A$ is bijective with purely
inseparable residue field extensions. Then every $b\in B$ has $b^{p^n}\in A$ for some $n$.
\end{theorem}

\begin{proof}
More generally let $A\subseteq B$ be radicial integral with $\operatorname{char}A=p$. The set $B'$ of $b\in B$ with $b^{p^n}\in A$ for some $n$ is an $A$-subalgebra
(additivity of Frobenius). If $B'\ne B$, Lemma~\ref{lem:gl-radicial-structure} applied to the radicial integral extension $B'\subseteq B$ gives $b\notin B'$ with
$b^2,b^3\in B'$, so $b^{p}\in B'$ and $b\in B'$; or a prime $\ell$ with $\ell b,b^\ell\in B'$, where $\ell=p$ gives $b^p\in B'$ and $\ell\ne p$ gives $b\in B'$ since $\ell$
is a unit. Both contradict $b\notin B'$. Stacks \texttt{0CNF}; the step that $B'\subseteq B$ is again radicial integral is left implicit there.
\end{proof}

\begin{remark}\label{rmk:gl-frobenius-descent}
\uses{thm:gl-frobenius-descent}
\cite{FHX} (proof of Theorem~8.3, positive characteristic): ``Since $\nu$ is a universal homeomorphism (Lemma~8.1), each generator has a $p$-power in $A_U$ by
[Stacks, Tag~0CNF]''. Stacks Lemma~29.47.9 (\texttt{0CNF}): for a ring map $A\to B$ inducing an isomorphism $A[1/p]\to B[1/p]$ (automatic in characteristic $p$),
$\Spec B\to\Spec A$ is a universal homeomorphism if and only if the kernel is locally nilpotent and every $b\in B$ has a $p$-power multiple and a $p$-power power in
the image of $A$; an integral, radicial, surjective morphism is a universal homeomorphism (Stacks \texttt{04DF}, \texttt{01S4}). Its proof uses the structure lemma
Stacks \texttt{0CNC} (an element outside a proper intermediate subalgebra with $b^2,b^3$ inside, or $\ell b,b^\ell$ inside for a prime $\ell$). They are formalized at ring level, not
imported; in characteristic $p$ no universal homeomorphisms or scheme theory are needed, the finite generation hypothesis is not used, and the conclusion is
Mathlib's \texttt{IsPRadical}.
\end{remark}

\begin{definition}[Subintegral towers]\label{def:gl-subintegral-tower}
\lean{EnvelopeProperty.IsSubintegralChain, EnvelopeProperty.Algebra.IsSubintegralTower, EnvelopeProperty.IsSubintegralChain.append, EnvelopeProperty.Algebra.isSubintegralTower_iff_forall_mem}
\leanok
\source{roberts-singh93-subintegrality:page-0010, stacks-universal-homeomorphisms:tag-0CND}
A finite sequence $e_1,\dots,e_n\in B$ is a \emph{subintegral chain} over $A$ if $e_i^2,e_i^3\in A[e_1,\dots,e_{i-1}]$ for every $i$. The extension $A\to B$ is
a \emph{subintegral tower} if every finite subset of $B$ lies in $A[e_1,\dots,e_n]$ for some subintegral chain; equivalently every element does
(chains concatenate). This is the definition of a subintegral extension of \cite[\S4]{RobertsSingh} and the form of Stacks \texttt{0CND}(2) over $\Q$.
\end{definition}

\begin{lemma}[Radicial integral extensions of $\Q$-algebras are subintegral]\label{lem:gl-subintegral-tower}
\lean{EnvelopeProperty.Algebra.IsRadicialIntegral.exists_sq_cube_of_algebraRat, EnvelopeProperty.Algebra.IsRadicialIntegral.isSubintegralTower, EnvelopeProperty.isSubintegralTower_of_maxSpec_injective}
\leanok
\uses{lem:gl-radicial-structure, def:gl-subintegral-tower, def:gl-radicial-integral}
\source{stacks-universal-homeomorphisms:tag-0CND}
A radicial integral extension $A\subseteq B$ of $\Q$-algebras is a subintegral tower. In particular, for $\operatorname{char}k=0$, a finite extension $A\subseteq B$ of
finitely generated $k$-algebras with $\operatorname{MaxSpec}B\to\operatorname{MaxSpec}A$ injective is a subintegral tower.
\end{lemma}

\begin{proof}
Over $\Q$ every prime $\ell$ is a unit, so Lemma~\ref{lem:gl-radicial-structure} gives $b\notin A$ with $b^2,b^3\in A$ whenever $A\ne B$ (Stacks \texttt{0CNB}). The
union $B'$ of all $A[e_1,\dots,e_n]$ over subintegral chains is a subalgebra, since chains concatenate. If $B'\ne B$, the radicial integral extension $B'\subseteq B$
contains $b\notin B'$ with $b^2,b^3\in B'$; then $b^2,b^3$ lie in some $A[e]$ for a chain $e$, and $(e,b)$ is a chain, so $b\in B'$, a contradiction. This replaces
the transfinite recursion of Stacks \texttt{0CND}.
\end{proof}

\begin{theorem}[Additive powers]\label{thm:gl-additive-powers}
\lean{EnvelopeProperty.Algebra.IsSubintegralTower.exists_addPow_mem_range, EnvelopeProperty.Algebra.IsSubintegralTower.isAddPowSubintegral, EnvelopeProperty.IsSubintegralChain.isAddPowSubintegral, EnvelopeProperty.IsAddPowSubintegral, EnvelopeProperty.addPowSubintegralClosure, EnvelopeProperty.IsAddPowSubintegral.add, EnvelopeProperty.IsAddPowSubintegral.mul, EnvelopeProperty.IsAddPowSubintegral.of_sq_of_cube, EnvelopeProperty.IsAddPowSubintegral.of_mem_adjoin_insert, EnvelopeProperty.IsAddPowSubintegral.exists_le, EnvelopeProperty.isAddPowSubintegral_of_tower, EnvelopeProperty.AddPow.hasExactSystem_of_eventually}
\leanok
\uses{lem:gl-radicial, def:gl-subintegral-tower}
\source{fhx-envelope-20261028:lem:normalization-local-sandwich, roberts-singh93-subintegrality:page-0020}
Let $A\subseteq B$ be a subintegral extension of $\Q$-algebras (a directed union of finite towers $A=A_0\subseteq A_1\subseteq\dots\subseteq A_n$
with $A_{i+1}=A_i[e_i]$ and $e_i^2,e_i^3\in A_i$; \cite[\S4]{RobertsSingh}, \cite[Tag~0CND]{Stacks}) and $b\in B$. There are integers $N\ge d\ge0$ and $c_1,\dots,c_d\in B$ such that
$a_m=b^m+\sum_{i=1}^d\binom mic_ib^{m-i}\in A$ for every $m\ge N$.
\end{theorem}

\begin{proof}
Call $b$ \emph{subintegral over $A$} if such $N$, $d$, $c_i$ exist. These elements form a subalgebra containing $A$, and $b$ is subintegral as soon as
$b^2$ and $b^3$ are (Roberts--Singh (4.8), (4.13)); by induction along a chain (Definition~\ref{def:gl-subintegral-tower}) every element of $B$ is then
subintegral, and $N\ge d$ after enlarging $N$. The Lean proof does not use Roberts--Singh's route through their Theorem~3.8. It turns a relation valid
only for $m\ge R$ into one valid for all $m\ge0$ with exponent $0$ by multiplying the $m$-th relation by $\pi(m)$, $\pi(X)=\prod_{j=1}^{R-1}(1-X/j)$
(``cut-off''). Sums use the product $C(T)e^{bT}$ of exponential generating functions, products the cut-off applied to $\binom mk\binom ml$, and the
square-and-cube step Lagrange interpolation over $\Q$ of $\binom ti\binom{(r-2t)/3}j$ along $2n+3m=r$ followed by the cut-off.
\end{proof}

\begin{remark}\label{rmk:gl-additive-powers}
\uses{thm:gl-additive-powers}
This is the first assertion of \cite[Lemma~8.2]{FHX}, which cites \cite[Theorem~4.17, Proposition~4.2, Corollary~4.3]{RobertsSingh} ; the first assertion is exactly that every
element of $B$ is subintegral over $A$ in the sense of \cite[p.~258]{RobertsSingh}, i.e.\ the implication (ii)$\Rightarrow$(i) of
\cite[Theorem~4.17]{RobertsSingh}. \cite[\S8]{FHX} calls
$A\subseteq B$ subintegral if it is integral, $\Spec B\to\Spec A$ is bijective and the residue field maps are isomorphisms, and states its equivalence
with the tower form (every finite subset of $B$ lies in a finite tower of steps $A'[e]$ with $e^2,e^3\in A'$, Stacks \texttt{0CND}); revision
\texttt{832e086} does not define the term. The tower form is the one Roberts--Singh use and the one Stacks \texttt{0CND} produces from a universal
homeomorphism in characteristic $0$, which is how \cite[Theorem~8.3]{FHX} obtains it. For an elementary step
$B=A[e]$ with $e^2,e^3\in A$ and $b=\alpha+\beta e$ one can take $d=1$, $c_1=-\beta e$: all terms of $b^m$ and of $m\beta eb^{m-1}$ except $m\alpha^{m-1}\beta e$ involve $e^i$,
$i\ge2$. A finite subintegral extension is a finite tower of such steps (Stacks \texttt{0CND}), but composing additive powers along a tower is the content of the cited
results; Lean follows the elementary route described in the proof above, and the theorem is not an explicit hypothesis. With Lemma~\ref{lem:gl-subintegral-tower} it gives the characteristic-$0$ input of \cite[Theorem~8.3]{FHX}
without imports.
\end{remark}

\begin{lemma}[Lagrange interpolation]\label{lem:gl-lagrange}
\lean{EnvelopeProperty.pow_mem_span_addPow, EnvelopeProperty.addPow_mem_span_pow, EnvelopeProperty.AddPow.eval_eq_sum_lagrange}
\leanok
\source{fhx-envelope-20261028:lem:normalization-local-sandwich}
Let $B$ be a $\Q$-algebra, $b\in B$, $d\ge0$ and $c_1,\dots,c_d\in B$, and put $a_m=b^m+\sum_{i=1}^{\min\{d,m\}}\binom mic_ib^{m-i}$ for $m\ge0$. Then
$b^{n+2d}\in(a_{n+d},\dots,a_{n+2d})B$ for every $n\ge0$, and $a_m\in b^{m-d}B$ for $m\ge d$. This is the second assertion of
\cite[Lemma~8.2]{FHX} (revision \texttt{832e086} states it for $n\ge N-d$ only).
\end{lemma}

\begin{proof}
Fix $m=n+d\ge d$, $c_0=1$ and $P(T)=\sum_{i=0}^d\binom Tic_ib^{m+d-i}\in B[T]$, of degree $\le d$. Then $P(m+r)=b^{d-r}a_{m+r}$ for $0\le r\le d$ and $P(0)=b^{m+d}$.
Lagrange interpolation over $\Q$ writes $P(0)$ as a $\Q$-combination of $P(m),\dots,P(m+d)$. The last claim is visible from the formula.
\end{proof}

\begin{theorem}[Psh functions descend through the normalization]\label{thm:gl-normalization-reflection}
\lean{EnvelopeProperty.MultiModel.IsUnibranch.isPSH_comp_symm_of_charP, EnvelopeProperty.MultiModel.IsUnibranch.exists_pow_pow_mem_sections, EnvelopeProperty.MultiModel.IsUnibranch.isFS_comp_symm_of_charP, EnvelopeProperty.MultiModel.IsUnibranch.isPSH_comp_symm, EnvelopeProperty.MultiModel.IsUnibranch.isPSH_comp_iff, EnvelopeProperty.MultiModel.IsUnibranch.hasEnvelopeProperty_of_normalization, EnvelopeProperty.MultiModel.IsUnibranch.hasEnvelopeProperty_normalization_iff}
\leanok
\uses{lem:gl-unibranch-universal, thm:gl-frobenius-descent, thm:gl-additive-powers, lem:gl-lagrange, lem:gl-sections, lem:gl-cones, thm:gl-fs-system-theory, def:gl-psh}
\source{fhx-envelope-20261028:thm:normalization-reflection, fhx-envelope-20261028:eq:normalization-global-sandwich}
Let $\mathcal X$ be unibranch, $\mathcal Y=\mathcal X^\nu$, and identify $\mathcal Y^\an$ with $\mathcal X^\an$ by $\nu^\an$. In characteristic $p$ (assuming
Theorem~\ref{thm:gl-frobenius-descent}), $\PSH(\mathcal Y,\theta)\subseteq\PSH(\mathcal X,\theta)$ for every $\theta\in\R^r$; in characteristic $0$ (assuming
Theorem~\ref{thm:gl-additive-powers}), for every $\theta\in N(\mathcal X)$, and for every $\theta\in\R^r$ by Theorem~\ref{thm:nn-normalization-reflection}
(one Lean statement, for every $\theta$). Both Theorems~\ref{thm:gl-frobenius-descent} and~\ref{thm:gl-additive-powers} are proved, so
there is no hypothesis.
\end{theorem}

\begin{proof}
\emph{Characteristic $p$.} $S^\nu$ is generated over $S$ by finitely many homogeneous $b_\ell$; on each chart $b_\ell/\tilde z_j^{\deg b_\ell}\in B_j$ has a $p$-power in $A_j$
(Theorem~\ref{thm:gl-frobenius-descent} and Lemma~\ref{lem:gl-unibranch-universal}); take one power $q$ for all $\ell,j$. Then $b_\ell^q\tilde z_j^{N\mathbf1}\in S$ for all $j$ and
$N\gg0$, so $b_\ell^q\in R(\mathcal X)$ (the monomials $z_j^{N\mathbf1}$ generate an ideal containing $\mathrm{Mon}(b')$, $b'\gg0$). Since Frobenius is additive and
$R(\mathcal X)$ is a ring, $p^q\in R(\mathcal X)$ for every $p\in S^\nu$, and for a section $f\in R_{ma}(\mathcal Y)$ the elements $f^q\tilde z^{q\beta}$ ($\beta\in\mathrm{Mon}(b)$) lie
in $R(\mathcal X)$, whence $f^q\in R_{qma}(\mathcal X)$ by the same ideal argument. If $\psi_\alpha=m^{-1}\max_t\{\log|f_t|+\lambda_t\}\in\mathcal H(\mathcal Y,a_\alpha)$
decrease to $\psi$, then $\psi_\alpha=(qm)^{-1}\max_t\{\log|f_t^q|+q\lambda_t\}\in\mathcal H(\mathcal X,a_\alpha)$, so $\psi\in\PSH(\mathcal X,\theta)$ by definition.

\emph{Characteristic $0$.} Claim: for $a\in\mathbf Z^r$ and $0\ne s\in R_a(\mathcal Y)$, $u=\log|s|\in\PSH(\mathcal X,a)$. For each chart $j$ let $f_j=s/\tilde z_j^a\in B_j$
(Lemma~\ref{lem:gl-sections}(4)); $A_j\subseteq B_j$ is subintegral (Lemma~\ref{lem:gl-unibranch-universal}); Theorem~\ref{thm:gl-additive-powers} gives $c_{ji}$, $a_{jm}$,
with $d,N$ common to the finitely many $j$. Fix $R$ such that (i) $\hat c_{ji}\coloneqq c_{ji}\tilde z_j^{R\mathbf1+ia}\in R_{ia+R\mathbf1}(\mathcal Y)$ for all $i\le d$; (ii) $\tilde
z_j^{R\mathbf1}$ kills the $z_j$-power torsion of the finite $S$-module $S^\nu/S$ (Noetherianity); (iii) $R\mathbf1+ia\ge0$ for $0\le i\le d$. For $m\ge N$,
$\tilde a_{jm}=\tilde z_j^{R\mathbf1}s^m+\sum_i\binom mi\hat c_{ji}s^{m-i}\in R_{ma+R\mathbf1}(\mathcal Y)$ satisfies $\tilde a_{jm}/\tilde z_j^{ma+R\mathbf1}=a_{jm}\in A_j$,
so by (ii) $t_{jm}\coloneqq\tilde z_j^{R\mathbf1}\tilde a_{jm}\in R_{ma+2R\mathbf1}(\mathcal X)$. For $n\ge\max\{1,N-d\}$ let $\mathcal S_n$ be the set of products of
$t_{jm}$ ($n+d\le m\le n+2d$) with monomials of degree $R\mathbf1+(n+2d-m)a$ (by (iii)), sections of $\mathcal X$ of degree $(n+2d)a+3R\mathbf1$. For $v\in\mathcal X^\an$:
every $t\in\mathcal S_n$ is $s^n$ times a section of $\mathcal Y$, so $\log|t|(v)\le n\,u(v)$; and for a chart with $|z^{(c)}_{j_c}|_v=1$,
Lemma~\ref{lem:gl-lagrange} gives $f_j^{n+2d}\in(a_{j,n+d},\dots,a_{j,n+2d})B_j$, and elements of $B_j$ have norm $\le1$ at $v$, so
$(n+2d)u(v)\le\max_m\log|a_{jm}|(v)\le\max_{t\in\mathcal S_n}\log|t|(v)$. At $\star$ the upper bound is $0$, so $\mathcal S_n$ contains a nonzero section, and
$F_n=(n+2d)^{-1}\max_{t\in\mathcal S_n}\log|t|\in\mathcal H_{\mathrm{hom}}(\mathcal X,a+\tfrac{3R}{n+2d}\mathbf 1)$ satisfies $u\le F_n\le\tfrac n{n+2d}u$. With $n_0>2d$,
$n_0\ge N-d$, $n_l=2^ln_0$ and $G_l=(1-2d/n_l)F_{n_l}$ one gets, as in \cite{FHX}, $G_{l+1}\le\tfrac{n_l-d}{n_l+d}u\le\tfrac{n_l-2d}{n_l}u\le G_l$; the $G_l$ are FS functions
decreasing to $u$ with classes tending to $a$ and $G_l(\star)=0$, so $u\in\PSH(\mathcal X,a)$ by definition.

For $\psi\in\PSH(\mathcal Y,\theta)$, $\theta\in N(\mathcal X)$, let $\psi_\alpha=m^{-1}\max_t\{\log|s_t|+\lambda_t\}\in\mathcal H(\mathcal Y,a_\alpha)$ decrease to $\psi$. By the claim
and Theorem~\ref{thm:gl-fs-system-theory}(1) (maxima and translates), $\psi_\alpha\in\PSH(\mathcal X,a_\alpha)$; PSH2 at $\theta\in N(\mathcal X)$
(Theorem~\ref{thm:gl-fs-system-theory}(2)) gives $\psi\in\PSH(\mathcal X,\theta)$.
\end{proof}

\begin{theorem}[Main theorem]\label{thm:gl-main}
\lean{EnvelopeProperty.MultiModel.hasEnvelopeProperty_of_isUnibranch, EnvelopeProperty.MultiModel.isPSH_uscReg_iSup_of_isUnibranch}
\leanok
\uses{lem:gl-unibranch-universal, cor:gl-normal-ep, thm:gl-normalization-reflection, lem:gl-pullback, lem:gl-cones}
\source{fhx-envelope-20261028:thm:intro-main}
Assume the hypotheses of Theorem~\ref{thm:gl-normal-ample-ep} (over the countable algebraically closed subfields of $k$; Theorems~\ref{thm:gl-frobenius-descent} and~\ref{thm:gl-additive-powers} are proved). Let $\mathcal X$ be a unibranch model (for instance the quotient model of an integral unibranch
$X\subseteq\PP^{M_1}\times\dots\times\PP^{M_r}$) and $\theta\in\R^r$ (Theorem~\ref{thm:nn-main}). Every nonempty family $(\varphi_i)\subseteq\PSH(\mathcal X,\theta)$ bounded above
satisfies $(\sup_i\varphi_i)^\star\in\PSH(\mathcal X,\theta)$.
\end{theorem}

\begin{proof}
Let $\mathcal Y=\mathcal X^\nu$, a normal model with $C(\mathcal X)\subseteq C(\mathcal Y)$ (sections pull back with the same values), so $\theta\in N(\mathcal Y)$. Identify
$\mathcal Y^\an=\mathcal X^\an$ by the homeomorphism $\nu^\an$ (Lemma~\ref{lem:gl-unibranch-universal}), which commutes with suprema and usc regularization. The
$\varphi_i\in\PSH(\mathcal Y,\theta)$ (Lemma~\ref{lem:gl-pullback}), so $(\sup\varphi_i)^\star\in\PSH(\mathcal Y,\theta)$ (Corollary~\ref{cor:gl-normal-ep}), hence
$\in\PSH(\mathcal X,\theta)$ (Theorem~\ref{thm:gl-normalization-reflection}).
\end{proof}

\begin{proposition}[Descent through the normalization under the envelope property]\label{prop:gl-conductor-descent}
\lean{EnvelopeProperty.MultiModel.HasEnvelopeProperty.exists_isPSH_comp_anMap_eq, EnvelopeProperty.MultiModel.exists_isConductorElem, EnvelopeProperty.MultiModel.IsConductorElem, EnvelopeProperty.MultiModel.IsConductorElem.eval_twistData, EnvelopeProperty.MultiModel.IsConductorElem.eq_of_anMap_eq_of_secAbs_ne_zero, EnvelopeProperty.MultiModel.HasEnvelopeProperty.anMap_ofNormalization_injective}
\leanok
\uses{lem:gl-pullback, lem:gl-usc-local, thm:gl-comparison, thm:gl-fs-system-theory, lem:gl-ep-basic, lem:gl-div-generic, rmk:gl-model-examples, lem:gl-sections}
\source{bj22-trivval:lem:descentpsh2, bj22-trivval:thm:descentpsh}
Let $\mathcal X$ be a model, $\mathcal Y=\mathcal X^\nu$, $\omega\in\R^r$ positive with $\mathrm{EP}_{\mathcal X}(\omega)$, and $\psi\in\PSH(\mathcal Y,\omega)$. There is
$\varphi\in\PSH(\mathcal X,\omega)$ with $\varphi\circ\nu^\an=\psi$.
\end{proposition}

\begin{proof}
$S^\nu$ is a finite $S$-module in the same fraction field, so the conductor $\{c\in S:cS^\nu\subseteq S\}$ contains a nonzero homogeneous $\tau$ of some degree $e\in\N^r$;
then $\tau R_b(\mathcal Y)\subseteq R_{b+e}(\mathcal X)$. Put $\ell=\log|\tau|\le0$ and $\Omega=\{\ell>-\infty\}\subseteq\mathcal X^\an$. Since $S[\tau^{-1}]=S^\nu[\tau^{-1}]$, $\nu^\an$
restricts to a homeomorphism $\nu^{-1}(\Omega)\to\Omega$. Let $\psi_\alpha\in\mathcal H(\mathcal Y,a_\alpha)$ decrease to $\psi$. For $\epsilon=1/M$, writing
$\psi_\alpha=(mM)^{-1}\max_t\log|s_t^M|+\dots$, the functions $\psi_\alpha+\epsilon\,\ell\circ\nu$ are pullbacks of FS functions $\varphi_{\alpha,\epsilon}\in\mathcal H(\mathcal X,
a_\alpha+\epsilon e)$ (use $s_t^M\tau^m$), decreasing on $\mathcal X^\an$ because $\nu^\an$ is surjective; so $\varphi_\epsilon=\lim_\alpha\varphi_{\alpha,\epsilon}\in\PSH(\mathcal
X,\omega+\epsilon e)$ and $\varphi_\epsilon\circ\nu=\psi+\epsilon\ell\circ\nu$. Fix $\delta>0$; for $\epsilon$ small, $\omega+\epsilon e\le(1+\delta)\omega$, the $\varphi_\epsilon$ increase as
$\epsilon\downarrow0$ and are bounded above, and $\mathrm{EP}((1+\delta)\omega)$ (Lemma~\ref{lem:gl-ep-basic}(2)) gives $\varphi^{(\delta)}=(\sup_\epsilon\varphi_\epsilon)^\star\in
\PSH((1+\delta)\omega)$. On $\Omega$, $\sup_\epsilon\varphi_\epsilon=\psi\circ(\nu|)^{-1}$, which is usc there, so $\varphi^{(\delta)}=\psi\circ(\nu|)^{-1}$ on $\Omega$
(Lemma~\ref{lem:gl-usc-local}(1)); as $\Omega\supseteq\mathcal X^{\divv}$, $\varphi^{(\delta)}$ is independent of $\delta$, and PSH2 at $\omega$ gives
$\varphi\in\PSH(\mathcal X,\omega)$. $\varphi\circ\nu$ and $\psi$ are $\omega$-psh on $\mathcal Y$ and agree on $\nu^{-1}(\Omega)\supseteq\mathcal Y^{\divv}$, hence
everywhere.
\end{proof}

\begin{corollary}[Characterization of unibranch varieties]\label{cor:gl-characterization}
\lean{EnvelopeProperty.MultiModel.isUnibranch_tfae_hasEnvelopeProperty, EnvelopeProperty.MultiModel.isUnibranch_of_exists_hasEnvelopeProperty, EnvelopeProperty.MultiModel.HasEnvelopeProperty.isUnibranch}
\leanok
\uses{thm:gl-main, prop:gl-conductor-descent, lem:gl-pl, def:gl-unibranch, lem:gl-pullback, lem:gl-charts}
\source{fhx-envelope-20261028:cor:intro-characterization, bj22-trivval:thm:unibranch}
Under the hypotheses of Theorem~\ref{thm:gl-main}, for a model $\mathcal X$ the following are equivalent: (1) $\mathcal X$ is unibranch; (2)
$\mathrm{EP}_{\mathcal X}(\theta)$ for every $\theta\in\R^r$ (Corollary~\ref{cor:nn-characterization}); (3) $\mathrm{EP}_{\mathcal X}(\omega)$ for some positive $\omega$. The implication
(3)$\Rightarrow$(1) uses none of the imported statements.
\end{corollary}

\begin{proof}
(1)$\Rightarrow$(2) is Theorem~\ref{thm:gl-main}; (2)$\Rightarrow$(3) take $\omega=\mathbf1$. (3)$\Rightarrow$(1) (\cite[Theorem~5.12]{BJ22}): FS functions on $\mathcal Y=\mathcal
X^\nu$ of classes $a\le\omega$ separate the points of $\mathcal Y^\an$ (two multigraded points differ on a homogeneous $p\in S^\nu_b$, and $n^{-1}\max\{\log|p|,-C\}\in
\mathcal H(\mathcal Y,b/n)$ with $b/n\le\omega$), and they are $\omega$-psh, hence of the form $\varphi\circ\nu^\an$ (Proposition~\ref{prop:gl-conductor-descent}). So
$\nu^\an$ is injective. Two maximal ideals of $B_j$ over one maximal ideal of $A_j$ give two rigid points of $U_j\subseteq\mathcal Y^\an$ ($|b|=1$ if $b\notin\mathfrak m$, $0$
otherwise, extended by Lemma~\ref{lem:gl-charts}) with the same image; so they coincide.
\end{proof}

\begin{remark}[Classes outside the model-nef cone]\label{rmk:gl-nonnef}
\uses{thm:gl-fs-system-theory, thm:gl-homogeneous-siu, prop:gl-homogeneous-slices, thm:gl-normalization-reflection, rmk:bj-psh-diagonal-finite}
PSH2 (\cite[Theorem~4.5]{BJ22}) is proved in \cite{BJ22} through \cite[Lemma~4.6]{BJ22} for the families $\mathcal H^{\mathrm{gf}}_\R(L)$, $L\in\Pic(X)_\Q$. Its proof has a gap for
families taking the value $-\infty$ (Remark~\ref{rmk:bj-psh-diagonal-finite}). The repairs used in this project need the family to be stable under $\max(\cdot,c)$ or to be
approximable by finite-valued members; both hold when $L$ is semiample and fail otherwise (constants are Fubini--Study functions of $L$ only if a multiple of $L$ is
globally generated). For $\theta\in N(\mathcal X)$ approximating classes can be raised into the semiample cone (Lemma~\ref{lem:gl-cones}(3)); for a psef class outside
$N(\mathcal X)$ they cannot. PSH2 at such a class is used in revision \texttt{832e086} of \cite{FHX} in the last step of the proof of Theorem~7.2, for the
limit $U$ in Proposition~7.4 (Proposition~7.5 of the revision of 2026-10-28), and in the characteristic-zero part of Theorem~8.3 (the definition-level limits in
Lemma~7.1, \cite[Corollary~6.17]{BJ22} and the characteristic-$p$ part of Theorem~8.3 are limits of Fubini--Study functions and need no PSH2); the revision of
2026-10-28 replaces these uses by \cite[Lemmas~2.3 and~7.4]{FHX} and by limits of Fubini--Study functions (Remark~\ref{rmk:nn-sources}). The version of Theorem~\ref{thm:gl-main} for all $\theta\in\R^r$ is obtained, with the same proofs, under the explicit
hypothesis
\begin{center}(\textsc{PSH2}$_{\mathcal X}$) \emph{for every $\theta\in\R^r$, a decreasing net $u_i\in\PSH(\mathcal X,\theta_i)$, $\theta_i\to\theta$, with $u=\lim u_i\not\equiv-\infty$ has
$u\in\PSH(\mathcal X,\theta)$,}\end{center}
for $\mathcal X$ and $\mathcal X^\nu$; this is \cite[Theorem~4.5]{BJ22} in the model. The Main Theorem covers all
classes. Chapter~\ref{sec:nonnef} shows that no hypothesis is needed: the three uses of PSH2 at non-nef classes in revision \texttt{832e086} are
repaired by homogeneous PSH2 and two structured diagonal lemmas, and the Main Theorem and Corollary~1.2 hold for every $\theta\in\R^r$
(Theorem~\ref{thm:nn-main}); the restriction $\theta\in N(\mathcal X)$ will be dropped from Theorems~\ref{thm:gl-homogeneous-siu},
\ref{thm:gl-normalization-reflection}, \ref{thm:gl-main}, Proposition~\ref{prop:gl-homogeneous-slices} and Corollaries~\ref{cor:gl-normal-ep},
\ref{cor:gl-characterization} when they are formalized.
\end{remark}

\begin{remark}[Trust boundary]\label{rmk:gl-trust-boundary}
\uses{thm:gl-main, cor:gl-characterization, thm:nn-main, cor:nn-characterization, thm:gl-dejong, thm:gl-fibre-constancy, thm:pg-psh-descent, thm:gl-smooth-ep, thm:import-sg-ep-char0, thm:kn-smooth-envelope-char-zero, thm:fu-variety-main, thm:fu-quot-main}
Theorems~\ref{thm:gl-main} and~\ref{thm:nn-main} (every $\theta\in\R^r$) and Corollaries~\ref{cor:gl-characterization} and~\ref{cor:nn-characterization}
use the proved alteration theorem above. The global compatibility statements retain three explicit hypotheses, each assumed only at the countable algebraically closed subfields of $k$: Lemma~5.5 of \cite{FHX}
(\texttt{LocalPositivity}, which gives Wu's regularization and Proposition~5.22 of \cite{FHX} in the form used by Corollary~\ref{thm:pg-psh-descent},
Chapter~\ref{sec:nf-route}); fibre constancy (Theorem~\ref{thm:gl-fibre-constancy}, supplied internally in the root public statements by
the proved fibre-constancy result in Chapter~\ref{sec:fc}); and Theorem~\ref{thm:import-sg-ep-char0} in characteristic~$0$ only. In characteristic~$p$
the last one is vacuous; in characteristic~$0$ it holds by Theorem~\ref{thm:kn-smooth-envelope-char-zero}, from relative ample Nadel
vanishing (Theorem~\ref{thm:kn-relative-ample-nadel-vanishing}). The implications (2)$\Rightarrow$(1) and (3)$\Rightarrow$(1) of the characterizations use none of them. Frobenius descent and additive
powers (Theorems~\ref{thm:gl-frobenius-descent}, \ref{thm:gl-additive-powers}) and PSH2 at non-nef classes (Chapter~\ref{sec:nonnef}) are proved, not assumed.
Nothing uses \cite[Theorems~5.6, 6.21, Lemma~4.6, Corollaries~4.17, 12.18]{BJ22}, test configurations, or the energy theory of \cite{BJ22}. The comparisons of
the model with arbitrary line bundles and $N^1$ (Remark~\ref{rmk:gl-model-comparison}) remain separate from the proved scheme-to-model alteration comparison in Remark~\ref{rmk:gl-dejong};
the statements for an integral projective variety and all of $N^1(X)_\R$ are Theorem~\ref{thm:var-main} and Corollary~\ref{cor:var-characterization}.
Lemma~5.5 of \cite{FHX} is proved over every algebraically closed field (Theorem~\ref{thm:lpn-local-positivity-instance}); its proof uses test
configurations of $X\times\mathbb A^1$, and Theorem~\ref{thm:lpn-main-unconditional} removes it from the hypotheses.
With both inputs proved, Theorems~\ref{thm:fu-variety-main} and~\ref{thm:fu-quot-main} and Corollaries~\ref{cor:fu-variety-characterization}
and~\ref{cor:fu-quot-characterization} state the main results of \cite{FHX} without hypotheses.
\end{remark}
